% #399 | Archive/2024/San pham - To Latex Strong - Tex hoa Lop 12 moi - Dot 2 (Chuong 1 + 2 + 3)/San pham - To Latex Strong - Tex hoa Lop 12 moi - Dot 2 (Chuong 1 + 2 + 3)/Data/Chuong 3/Ket noi tri thuc/To 22 - Dot 16 - Chuong 3 - Bai 2 - KNTT.tex | mã gốc 2D3H2-2
\begin{ex}%[2D3H2-2]%[Tổ 16 - Đợt 16 - Chương 3 - Bài 2 - KNTT]%[Sauluoi3105]
Một trang trại nuôi gà có cuộc khảo sát về khối lượng (đơn vị gam) của $400$ quả trứng, số liệu được ghi lại trong bảng sau đây
\begin{center}
\begin{tabular}{|c|c|c|}
\hline Lớp khối lượng $($ gam $)$ & Giá trị đại diện & Tần số \\
\hline$[27{,}5 ; 32{,}5)$ & 30 & 18 \\
\hline$[32{,}5 ; 37{,}5)$ & 35 & 76 \\
\hline$[37{,}5 ; 42{,}5)$ & 40 & 200 \\
\hline$[42{,}5 ; 47{,}5)$ & 45 & 100 \\
\hline$[47{,}5 ; 52{,}5)$ & 50 & 6 \\
\hline & & $n=400$. \\
\hline
\end{tabular}
\end{center}
Tính phương sai của mẫu số liệu.
\loigiai{
Ta có $\bar{x}=\dfrac{m_1 x_1+m_2 x_2+\ldots+m_5 x_5}{n} \approx 40$.\\
Phương sai $s^2=\dfrac{m_1\left(x_1-\bar{x}\right)^2+\ldots+m_5\left(x_5-\bar{x}\right)^2}{n} \approx 17$.
}
\end{ex}

% #400 | Archive/2024/San pham - To Latex Strong - Tex hoa Lop 12 moi - Dot 2 (Chuong 1 + 2 + 3)/San pham - To Latex Strong - Tex hoa Lop 12 moi - Dot 2 (Chuong 1 + 2 + 3)/Data/Chuong 3/Ket noi tri thuc/To 22 - Dot 16 - Chuong 3 - Bai 2 - KNTT.tex | mã gốc 2D3H2-2
\begin{ex}%[2D3H2-2]%[Tổ 16 - Đợt 16 - Chương 3 - Bài 2 - KNTT]%[Sauluoi3105]
Khối lượng (đơn vị kg ) của $20$ con cá được cho bởi bảng sau đây
\begin{center}
\begin{tabular}{|l|c|c|c|c|c|}
\hline Lớp khối lượng $(\mathrm{kg})$ & {$[0{,}6 ; 0{,}8)$} & {$[0{,}8 ; 1{,}0)$} & {$[1{,}0 ; 1{,}2)$} & {$[1{,}2 ; 1{,}4)$} & $\mathbf{n}$ \\
\hline Giá trị đại diện & 0{,}7 & 0{,}9 & 1{,}1 & 1{,}3 & \\
\hline Tần số & 4 & 6 & 6 & 4 & $20$ \\
\hline
\end{tabular}
\end{center}
Tính độ lệch chuẩn của mẫu số liệu.
\loigiai{
Ta có $\bar{x}=\dfrac{m_1 x_1+m_2 x_2+\cdots+m_4 x_4}{n} \approx 1$.\\
Phương sai $s^2=\dfrac{m_1\left(x_1-\bar{x}\right)^2+\cdots+m_4\left(x_4-\bar{x}\right)^2}{n} \approx 0{,}042$.\\
Độ lệch chuẩn $s=\sqrt{s^2} \approx 0{,}205$.
}
\end{ex}

% #401 | Archive/2024/San pham - To Latex Strong - Tex hoa Lop 12 moi - Dot 2 (Chuong 1 + 2 + 3)/San pham - To Latex Strong - Tex hoa Lop 12 moi - Dot 2 (Chuong 1 + 2 + 3)/Data/Chuong 3/Ket noi tri thuc/To 22 - Dot 16 - Chuong 3 - Bai 2 - KNTT.tex | mã gốc 2D3N2-2
\begin{ex}%[2D3N2-2]%[Tổ 16 - Đợt 16 - Chương 3 - Bài 2 - KNTT]%[Sauluoi3105]
Tìm hiểu thời gian chạy cự li $1000$ m ( đơn vị: giây) của các bạn học sinh trong một lớp thu được kết quả sau
\begin{center}
\begin{tabular}{|c|c|c|c|c|c|}
\hline Thời gian & {$[125 ; 127)$} & {$[127 ; 129)$} & {$[129 ; 131)$} & {$[131 ; 133)$} & {$[133 ; 135)$} \\
\hline Số bạn & 3 & 7 & 15 & 10 & 5 \\
\hline
\end{tabular}
\end{center}
Nhóm $[131 ; 133)$ có tần số là
\loigiai{
Nhóm $[131 ; 133)$ có tần số là $10$.
}
\end{ex}

% #402 | Archive/2024/San pham - To Latex Strong - Tex hoa Lop 12 moi - Dot 2 (Chuong 1 + 2 + 3)/San pham - To Latex Strong - Tex hoa Lop 12 moi - Dot 2 (Chuong 1 + 2 + 3)/Data/Chuong 3/Ket noi tri thuc/To 22 - Dot 16 - Chuong 3 - Bai 2 - KNTT.tex | mã gốc 2D3N2-2
\begin{ex}%[2D3N2-2]%[Tổ 16 - Đợt 16 - Chương 3 - Bài 2 - KNTT]%[Sauluoi3105]
Cho mẫu số liệu ghép nhóm về khối lượng (đơn vị: gram) của 30 củ khoai tây như sau:
\begin{center}
\begin{tabular}{|c|c|c|c|c|c|}
\hline Khối lượng & {$[70 ; 80)$} & {$[80 ; 90)$} & {$[90 ; 100)$} & {$[100 ; 110)$} & {$[110 ; 120)$} \\
\hline Tần số & 3 & 6 & 12 & 6 & 3 \\
\hline
\end{tabular}
\end{center}
Giá trị đại diện của nhóm $[90 ; 100)$.
\loigiai{
Giá trị đại diện của nhóm $[90 ; 100)$ là $\dfrac{90+100}{2}=95$.
}
\end{ex}

% #403 | Archive/2024/San pham - To Latex Strong - Tex hoa Lop 12 moi - Dot 2 (Chuong 1 + 2 + 3)/San pham - To Latex Strong - Tex hoa Lop 12 moi - Dot 2 (Chuong 1 + 2 + 3)/Data/Chuong 3/Ket noi tri thuc/To 22 - Dot 16 - Chuong 3 - Bai 2 - KNTT.tex | mã gốc 2D3N2-2
\begin{ex}%[2D3N2-2]%[Tổ 16 - Đợt 16 - Chương 3 - Bài 2 - KNTT]%[Sauluoi3105]
Trong bài thực hành đo hiệu điện thế của mạch điện, bạn Minh tiến hành đo $12$ lần, kết quả như sau:
\begin{center}
\begin{tabular}{|l|c|c|c|c|}
\hline
Hiệu điện thế đo được (Vôn) & {$[3{,}85 ; 3{,}90)$} & {$[3{,}90 ; 3{,}95)$} & {$[3{,}95 ; 4{,}00)$} & {$[4{,}00 ; 4{,}05)$} \\
\hline Số lần đo & 2 & 3 & 5 & 2 \\
\hline
\end{tabular}
\end{center}
Tính độ lệch chuẩn của mẫu số liệu ghép nhóm trên ( làm tròn đến chư số thập phân thứ ba ).
\loigiai{
Chọn giá trị đại diện cho nhóm số liệu ta có
\begin{center}
\begin{tabular}{|c|c|c|c|c|}
\hline Giá trị đại diện & 3{,}875 & 3{,}925 & 3{,}975 & 4{,}025 \\
\hline Số lần đo & 2 & 3 & 5 & 2 \\
\hline
\end{tabular}
\end{center}
Hiệu điện thế trung bình của mạch điện là
$$\bar{x}=\dfrac{1}{12}(2\cdot 3{,}875+3\cdot 3{,}925+5\cdot 3{,}975+2\cdot 4{,}025) \approx 3{,}954(\mathrm{~V}).$$
Độ lệch chuẩn của hiệu điện thế của mạch điện là
$$s=\sqrt{\dfrac{1}{12}\left(2\cdot 3{,}875^2+3\cdot 3{,}925^2+5\cdot 3{,}975^2+2\cdot 4{,}025^2\right)-3{,}954^2} \approx 0{,}05.$$
}
\end{ex}

% #404 | Archive/2024/San pham - To Latex Strong - Tex hoa Lop 12 moi - Dot 2 (Chuong 1 + 2 + 3)/San pham - To Latex Strong - Tex hoa Lop 12 moi - Dot 2 (Chuong 1 + 2 + 3)/Data/Chuong 3/Ket noi tri thuc/To 22 - Dot 16 - Chuong 3 - Bai 2 - KNTT.tex | mã gốc 2D3N2-2
\begin{ex}%[2D3N2-2]%[Tổ 22 - Đợt 16 - Chương 3 - Bài 2 - KNTT]%[Sauluoi3105]
Thời gian chạy tập luyện cự li $100 m$ của một vận động viên được cho trong bảng sau
\begin{center}
\begin{tabular}{|l|c|c|c|c|c|}
\hline Thời gian (giây) & {$[10 ; 10{,}4)$} & {$[10{,}4 ; 10{,}8)$} & {$[10{,}8 ; 11{,}2)$} & {$[11{,}2 ; 11{,}6)$} & {$[11{,}6 ; 12{,}0)$} \\
\hline Số lần chạy & 3 & 8 & 6 & 2 & 1 \\
\hline
\end{tabular}
\end{center}
Tính phương sai của mẫu số liệu ghép nhóm trên.
\loigiai{
Chọn giá trị đại diện cho nhóm số liệu ta có
\begin{center}
\begin{tabular}{|l|c|c|c|c|c|}
\hline Giá trị đại diện & 10{,}2 & 10{,}6 & 11 & 11{,}4 & 11{,}8 \\
\hline Số lần chạy & 3 & 8 & 6 & 2 & 1 \\
\hline
\end{tabular}
\end{center}
Thời gian trung bình mỗi lần chạy của vận động viên trên là
$$\bar{x}=\dfrac{1}{20}(3\cdot 10{,}2+8 \cdot 10{,}6+6\cdot 11+2\cdot 11{,}4+1\cdot 11{,}8)=10{,}8\,\text{(giây).}$$
Phương sai của thời gian chạy của vận động viên trên là
$$s^2=\dfrac{1}{20}\left(3\cdot 10{,}2^2+8\cdot 10{,}6^2+6\cdot 11^2+2\cdot 11{,}4^2+1\cdot 11{,}8^2\right)-10{,}8^2=0{,}17.$$
}
\end{ex}

% #1135 | Archive/2024/0. [LVH] BG - 12 New - 4in1-20240625T224301Z-001/0. [LVH] BG - 12 New - 4in1/2. [LVH] TEX SGK 12 NEW-20240418T074833Z-001/2. [LVH] TEX SGK 12 NEW/[LVH] SGK-12-KNTT/data/C3/12-SGK-KNTT-3-10.tex | mã gốc 2D3H2-2
\begin{vd}%[2D3H2-2]
	Người ta theo dõi sự thay đổi cân nặng, được tính bằng hiệu cân nặng trước và sau ba tháng áp dụng chế độ ăn kiêng của một số người cho kết quả như sau:
\begin{center}
		\begin{tabular}{|l|c|c|c|c|c|}
		\hline Thay đổi cân nặng (kg) & {$[-1; 0)$} & {$[0; 1)$} & {$[1; 2)$} & {$[2; 3)$} & {$[3; 4)$} \\
		\hline Số người nam & $2$ & $3$ & $5$ & $3$ & $2$ \\
		\hline Số người nữ & $2$ & $7$ & $12$ & $7$ & $2$ \\
		\hline
	\end{tabular}
\end{center}	
	Tính số trung bình, phương sai, độ lệch chuẩn và nhận xét về sự thay đổi cân nặng của người nam, người nữ sau ba tháng áp dụng chế độ ăn kiêng.
\loigiai{Chọn giá trị đại diện cho các nhóm số liệu, ta có	
\begin{center}
		\begin{tabular}{|l|c|c|c|c|c|}
		\hline Giá trị đại diện & $-0{,}5$ & $0{,}5$ & $1{,}5$ & $2{,}5$ & $3{,}5$ \\
		\hline Số người nam & $2$ & $3$ & $5$ & $3$ & $2$ \\
		\hline Số người nữ & $2$ & $7$ & $12$ & $7$ & $2$ \\
		\hline
	\end{tabular}
\end{center}	
	Tổng số người nam là $n_1=2+3+5+3+2=15$.\\
	Tổng số người nữ là $n_2=2+7+12+7+2=30$.\\
	Thay đổi cân nặng trung bình của người nam là:
	$$
	\overline{x}_1=\frac{1}{15}[2 \cdot(-0{,}5)+3 \cdot 0{,}5+5 \cdot 1{,}5+3 \cdot 2{,}5+2 \cdot 3{,}5]=1{,}5~(\mathrm{kg}).
	$$
	
	Thay đổi cân nặng trung bình của người nữ là:
	$$
	\overline{x}_2=\dfrac{1}{30}[2 \cdot(-0{,}5)+7 \cdot 0{,}5+12 \cdot 1{,}5+7 \cdot 2{,}5+2 \cdot 3{,}5]=1{,}5~(\mathrm{kg}) .
	$$	
	Phương sai và độ lệch chuẩn của mẫu số liệu về thay đổi cân nặng của người nam là
	$$
	s_1^2=\dfrac{1}{15}\left[2 \cdot(-0{,}5)^2+3 \cdot 0{,}5^2+5 \cdot 1{,}5^2+3 \cdot 2{,}5^2+2 \cdot 3{,}5^2\right]-1{,}5^2 \approx 1{,}21^2; s_1 \approx 1{,}21.
	$$
	
	Phương sai và độ lệch chuẩn của mẫu số liệu về thay đổi cân nặng của người nữ là
	$$
	s_2^2=\dfrac{1}{30}\left[2 \cdot(-0{,}5)^2+7 \cdot 0{,}5^2+12 \cdot 1{,}5^2+7 \cdot 2{,}5^2+2 \cdot 3{,}5^2\right]-1{,}5^2 \approx 2{,}06^2; s_2 \approx 2{,}06.
	$$	
	Như vậy, sau ba tháng áp dụng chế độ ăn kiêng này, về trung bình sự thay đổi cân nặng của nam và nữ là như nhau. Tuy nhiên, sự biến động về thay đổi cân nặng của nữ nhiều hơn so với của nam.		
}
	
\end{vd}

% #1136 | Archive/2024/0. [LVH] BG - 12 New - 4in1-20240625T224301Z-001/0. [LVH] BG - 12 New - 4in1/2. [LVH] TEX SGK 12 NEW-20240418T074833Z-001/2. [LVH] TEX SGK 12 NEW/[LVH] SGK-12-KNTT/data/C3/12-SGK-KNTT-3-10.tex | mã gốc 2D3H2-2
\begin{vd}%[2D3H2-2]
	Anh An đầu tư số tiền bằng nhau vào hai lĩnh vực kinh doanh $A, B$. Anh An thống kê số tiền thu được mỗi tháng trong vòng $60$ tháng theo mỗi lĩnh vực cho kết quả như sau:
	\begin{center}
		\begin{tabular}{|l|c|c|c|c|c|}
			\hline Số tiền (triệu đồng) & {$[5 ; 10)$} & {$[10 ; 15)$} & {$[15 ; 20)$} & {$[20 ; 25)$} & {$[25 ; 30)$} \\
			\hline Số tháng đầu tư vào lĩnh vực $A$ & $5$ & $10$ & $30$ & $10$ & $5$ \\
			\hline Số tháng đầu tư vào lĩnh vực $B$ & $20$ & $5$ & $10$ & $5$ & $20$ \\
			\hline
		\end{tabular}
	\end{center}	
	So sánh giá trị trung bình và độ lệch chuẩn của số tiền thu được mỗi tháng khi đầu tư vào mỗi lĩnh vực $A$, $B$. Đầu tư vào lĩnh vực nào \lq\lq rủi ro\rq\rq ~hơn?
\end{vd}

% #1137 | Archive/2024/0. [LVH] BG - 12 New - 4in1-20240625T224301Z-001/0. [LVH] BG - 12 New - 4in1/2. [LVH] TEX SGK 12 NEW-20240418T074833Z-001/2. [LVH] TEX SGK 12 NEW/[LVH] SGK-12-KNTT/data/C3/12-SGK-KNTT-3-10.tex | mã gốc 2D3H2-2
\begin{vd}%[2D3H2-2]
	Thống kê lợi nhuận hàng tháng (đơn vị: triệu đồng) trong $20$ tháng của hai nhà đầu tư được cho như sau:
	\begin{center}
		\begin{tabular}{|l|c|c|c|c|c|}
			\hline Lợi nhuận & {$[10; 20)$} & {$[20; 30)$} & {$[30; 40)$} & {$[40; 50)$} & {$[50 ; 60)$} \\
			\hline Số tháng & $2$ & $4$ & $8$ & $4$ & $2$ \\
			\hline
		\end{tabular}
	\end{center}	
	\begin{center}
		Bảng 3.2. Lợi nhuận theo tháng của nhà đầu tư nhỏ
	\end{center}
	\begin{center}
		\begin{tabular}{|l|c|c|c|c|c|}
			\hline Lợi nhuận & {$[510; 520)$} & {$[520; 530)$} & {$[530; 540)$} & {$[540; 550)$} & {$[550; 560)$} \\
			\hline Số tháng & $4$ & $3$ & $6$ & $3$ & $4$ \\
			\hline
		\end{tabular}\\
	Bảng 3.3. Lợi nhuận theo tháng của nhà đầu tư lớn
	\end{center}	
	Tính độ lệch chuẩn của hai mẫu số liệu ghép nhóm trên. Có nên dựa vào độ lệch chuẩn để so sánh độ rủi ro của hai nhà đầu tư này không?
\loigiai{
Chọn điểm đại diện cho các nhóm số liệu ta tính được các số đặc trưng như sau\\
Lợi nhuận trung bình một tháng của các nhà đầu tư tương ứng là
$$
\bar{x}_A=\dfrac{1}{20}(2 \cdot 15+\ldots+2 \cdot 55)=35 \text { (triệu đồng); } \bar{x}_B=\dfrac{1}{20}(4 \cdot 515+\ldots+4 \cdot 555)=535 \text { (triệu đồng). }
$$
Độ lệch chuẩn của lợi nhuận hàng tháng của hai nhà đầu tư tương ứng là
$$s_A=\sqrt{\dfrac 1{20}\left(2\cdot 15^2+\ldots+2\cdot 55^2\right)-(35)^2}\approx 10{,}95;$$
$$s_B=\sqrt{\dfrac 1{20}\left(4\cdot 515^2+\ldots+4\cdot 555^2\right)-(535)^2}\approx 13{,}78.$$
Độ lệch chuẩn cho lợi nhuận hàng tháng của nhà đầu tư lớn cao hơn của nhà đầu tư nhỏ. Lợi nhuận trung bình của hai nhà đầu tư khác nhau rất nhiều, do đó ta không nên dùng độ lệch chuẩn để so sánh mức độ rủi ro của hai nhà đầu tư này.
}
\end{vd}

% #1138 | Archive/2024/0. [LVH] BG - 12 New - 4in1-20240625T224301Z-001/0. [LVH] BG - 12 New - 4in1/2. [LVH] TEX SGK 12 NEW-20240418T074833Z-001/2. [LVH] TEX SGK 12 NEW/[LVH] SGK-12-KNTT/data/C3/12-SGK-KNTT-3-10.tex | mã gốc 2D3H2-2
\begin{bt}%[2D3H2-2]
	Kiểm tra khối lượng của $30$ bao xi măng (đơn vị: kg) được chọn ngẫu nhiên trước khi xuất xưởng cho kết quả như sau:
\begin{center}
		\begin{tabular}{llllllllll}
		$49{,}5$ & $51{,}1$ & $50{,}8$ & $50{,}2$ & $48{,}7$ & $49{,}6$ & $51{,}3$ & $51{,}4$ & $50{,}1$ & $50{,}5$ \\
		$48{,}9$ & $49{,}3$ & $50{,}7$ & $48{,}8$ & $49{,}8$ & $48{,}8$ & $51{,}2$ & $50{,}4$ & $50{,}0$ & $51{,}2$ \\
		$51{,}4$ & $48{,}7$ & $51{,}2$ & $50{,}6$ & $50{,}9$ & $49{,}2$ & $50{,}7$ & $51{,}1$ & $48{,}6$ & $49{,}6$.
	\end{tabular}
\end{center}	
	\begin{enumerate}[a)]
		\item Thay dấu \lq\lq $?$\rq\rq ~bằng số thích hợp để hoàn thiện mẫu số liệu ghép nhóm sau.
		\begin{center}
			\begin{tabular}{|l|c|c|c|c|c|c|}
				\hline Nhóm số liệu & {$[48{,}5; 49)$} & {$[49; 49{,}5)$} & {$[49{,}5; 50)$} & {$[50 ; 50{,}5)$} & {$[50{,}5 ; 51)$} & {$[51 ; 51{,}5)$} \\
				\hline Số bao xi măng & $?$ & $?$ & $?$ & $?$ & $?$ & $?$ \\
				\hline
			\end{tabular}
		\end{center}
	\item Tính phương sai và độ lệch chuẩn của mẫu số liệu gốc và mẫu số liệu ghép nhóm. Giá trị nào là giá trị chính xác? Giá trị nào là giá trị xấp xỉ?
	\end{enumerate}
\loigiai{
	\begin{enumerate}[a)]
		\item 
	\begin{center}
		\begin{tabular}{|l|c|c|c|c|c|c|}
			\hline Nhóm số liệu & {$[48{,}5; 49)$} & {$[49; 49{,}5)$} & {$[49{,}5; 50)$} & {$[50 ; 50{,}5)$} & {$[50{,}5; 51)$} & {$[51; 51{,}5)$} \\
			\hline Số bao xi măng & $6$ & $2$ & $4$ & $4$ & $6$ & $8$ \\
			\hline
		\end{tabular}
	\end{center}
	\item $\text{}$
	\begin{itemize}
		\item Phương sai và độ lệch chuẩn của mẫu số liệu gốc.\\		
			Khối lượng trung bình của $30$ bao xi măng là 
			$$\overline{x}=\dfrac{48{,}6+...+51{,}4}{30}=\dfrac{1504{,}3}{30}=50{,}1433.$$
			Phương sai $s_x^2=\dfrac{1}{30}$$\left[(48{,}6-50{,}1433)^2+...+(51{,}4-50{,}1433)^2\right]=0{,}871789$.\\
			Độ lệch chuẩn $s_x=\sqrt{s_x^2}=\sqrt{0{,}871789}=0{,}933696$.
		\item Phương sai và độ lệch chuẩn của mẫu số liệu ghép nhóm.
			\begin{center}
			\begin{tabular}{|l|c|c|c|c|c|c|}
				\hline Nhóm số liệu & {$[48,5; 49)$} & {$[49; 49,5)$} & {$[49,5; 50)$} & {$[50 ; 50,5)$} & {$[50,5 ; 51)$} & {$[51 ; 51,5)$} \\
				\hline Giá trị đại diện &$48{,}75$ &$49{,}25$ & $49{,}75$&$50{,}25$ &$50{,}75$ & $51{,}25$\\
				\hline Số bao xi măng & $6$ & $2$ & $4$ & $4$ & $6$ & $8$ \\
				\hline
			\end{tabular}
		\end{center}
		Khối lượng trung bình của $30$ bao xi măng là 	
		$$\overline{x}=\dfrac{6\cdot48{,}75+2\cdot49{,}25+4\cdot49{,}75+4\cdot20{,}25+6\cdot50{,}75+8\cdot51{,}25}{30}=50{,}18.$$
		Phương sai 		
		$$s^2=\dfrac{1}{30}\left(6\cdot48{,}75^2+2\cdot49{,}25^2+4\cdot49{,}75^2+4\cdot20{,}25^2+6\cdot50{,}75^2+8\cdot51{,}25^2\right)-50{,}18^2=1{,}197.$$
		Độ lệch chuẩn $s_x=\sqrt{s_x^2}=\sqrt{1{,}197}=1{,}094$.
	\end{itemize}	
	\end{enumerate}
Phương sai, độ lệch chuẩn của mẫu số liệu ghép nhóm là các xấp xỉ cho phương sai, độ lệch chuẩn của mẫu số liệu gốc.
}
\end{bt}

% #1139 | Archive/2024/0. [LVH] BG - 12 New - 4in1-20240625T224301Z-001/0. [LVH] BG - 12 New - 4in1/2. [LVH] TEX SGK 12 NEW-20240418T074833Z-001/2. [LVH] TEX SGK 12 NEW/[LVH] SGK-12-KNTT/data/C3/12-SGK-KNTT-3-10.tex | mã gốc 2D3H2-2
\begin{bt}%[2D3H2-2]
	Tuổi thọ của một số linh kiện điện tử (đơn vị: năm) được sản xuất bởi hai phân xưởng được cho như sau:
	\begin{center}
		\begin{tabular}{|l|c|c|c|c|c|}
			\hline Tuổi thọ (năm) & {$[1{,}5; 2)$} & {$[2; 2{,}5)$} & {$[2{,}5; 3)$} & {$[3; 3{,}5)$} & {$[3{,}5; 4)$} \\
			\hline Số linh kiện của phân xưởng $1$ & $4$ & $9$ & $13$ & $8$ & $6$ \\
			\hline Số linh kiện của phân xưởng $2$ & $2$ & $8$ & $20$ & $7$ & $3$ \\
			\hline
		\end{tabular}
	\end{center}	
	Tính phương sai và độ lệch chuẩn của mỗi mẫu số liệu ghép nhóm và nhận xét về độ phân tán của tuổi thọ các linh kiện điện tử được sản xuất bởi mỗi phân xưởng.
	\loigiai{
		\begin{itemize}
		\item 	Phương sai và độ lệch chuẩn của của số linh kiện ở phân xưởng $1$.
			\begin{center}
			\begin{tabular}{|l|c|c|c|c|c|}
				\hline Tuổi thọ (năm) & {$[1{,}5; 2)$} & {$[2; 2{,}5)$} & {$[2{,}5; 3)$} & {$[3; 3{,}5)$} & {$[3{,}5; 4)$} \\
				\hline Giá trị đại diện &$1{,}75$ & $2{,}25$ & $2{,}75$ & $3{,}25$ & $3{,}75$\\
				\hline Số linh kiện của phân xưởng $1$ & $4$ & $9$ & $13$ & $8$ & $6$ \\				
				\hline
			\end{tabular}
		\end{center}
	Tuổi thọ trung bình của $40$ linh kiện là
	$$\overline{x}_1=\dfrac{1}{40}(4\cdot 1{,}75+9\cdot 2{,}25+13\cdot 2{,}75+8\cdot 3{,}25+6\cdot 3{,}75)=2{,}7875.$$
		Phương sai và độ lệch chuẩn	của nhóm linh kiện ở phân xưởng $1$ là
		$$s_1^2=\dfrac{1}{40}\left(4\cdot 1{,}75^2+9\cdot 2{,}25^2+13\cdot 2{,}75^2+8\cdot 3{,}25^2+6\cdot 3{,}75^2\right)-2{,}7875=0{,}35484. $$
		Suy ra $s_1=\sqrt{0{,}35484}=0{,}59569.$
		\item 	Phương sai và độ lệch chuẩn của của số linh kiện ở phân xưởng $2$.
			\begin{center}
			\begin{tabular}{|l|c|c|c|c|c|}
				\hline Tuổi thọ (năm) & {$[1{,}5 ; 2)$} & {$[2 ; 2{,}5)$} & {$[2{,}5 ; 3)$} & {$[3 ; 3{,}5)$} & {$[3{,}5 ; 4)$} \\	
				\hline Giá trị đại diện &$1{,}75$ & $2{,}25$ & $2{,}75$ & $3{,}25$ & $3{,}75$\\		
				\hline Số linh kiện của phân xưởng $2$ & $2$ & $8$ & $20$ & $7$ & $3$ \\					
				\hline
			\end{tabular}
		\end{center}
			Tuổi thọ trung bình của $40$ linh kiện là
			$$\overline{x}_2=\dfrac{1}{40}(2\cdot 1{,}75+8\cdot 2{,}25+20\cdot 2{,}75+7\cdot 3{,}25+3\cdot 3{,}75)=2{,}7625.$$
			Phương sai và độ lệch chuẩn	của nhóm linh kiện ở phân xưởng $2$ là
			$$s_2^2=\dfrac{1}{40}\left(2\cdot 1{,}75^2+8\cdot 2{,}25^2+20\cdot 2{,}75^2+7\cdot 3,25^2+3\cdot 3{,}75^2\right)-2{,}7625^2=0{,}2186. $$
			Suy ra $s_2=\sqrt{0{,}2186}=0{,}4675.$	
		\end{itemize}
	}
\end{bt}

% #1140 | Archive/2024/0. [LVH] BG - 12 New - 4in1-20240625T224301Z-001/0. [LVH] BG - 12 New - 4in1/2. [LVH] TEX SGK 12 NEW-20240418T074833Z-001/2. [LVH] TEX SGK 12 NEW/[LVH] SGK-12-KNTT/data/C3/12-SGK-KNTT-3-10.tex | mã gốc 2D3H2-2
\begin{bt}%[2D3H2-2]
	Một nhóm $20$ học sinh dùng một thiết bị đo đường kính của một nhân tế bào cho kết quả như sau:
	\begin{center}
		\begin{tabular}{|l|c|c|c|c|}
			\hline Kết quả đo $(\mu \mathrm{m})$ & {$[4{,}5; 5)$} & {$[5; 5{,}5)$} & {$[5{,}5; 6)$} & {$[6; 6{,}5)$} \\
			\hline Số học sinh & $3$ & $8$ & $7$ & $2$ \\
			\hline
		\end{tabular}
	\end{center}	
	\begin{enumerate}[a)]
	\item 	Tính số trung bình và độ lệch chuẩn của mẫu số liệu ghép nhóm trên.
	\item Số trung bình và độ lệch chuẩn cho biết thông tin gì?
	\end{enumerate}
\loigiai{	
	\begin{enumerate}[a)]
		\item 	Ta có
			\begin{center}
			\begin{tabular}{|l|c|c|c|c|}
				\hline Kết quả đo $(\mu \mathrm{m})$ & {$[4{,}5 ; 5)$} & {$[5 ; 5{,}5)$} & {$[5{,}5 ; 6)$} & {$[6 ; 6{,}5)$} \\
				\hline Giá trị đại diện & $4{,}75$ & $5{,}25$ & $5{,}75$ & $6{,}25$ \\
				\hline Số học sinh & $3$ & $8$ & $7$ & $2$ \\
				\hline
			\end{tabular}		
		\end{center}
		Số đo trung bình 
		$$\overline{x}=\dfrac{1}{20}(3\cdot4{,}75+8\cdot5{,}25+7\cdot5{,}75+2\cdot6{,}25)=5{,}45.$$
		Độ lệ chuẩn
		$$s=\sqrt{\dfrac{1}{20}\left(3\cdot4{,}75^2+8\cdot5{,}25^2+7\cdot5{,}75^2+2\cdot6{,}25^2\right)-5{,}45^2}=0{,}43012.$$	
		\item Số trung bình và độ lệch chuẩn cho biết sự chênh lệch về giá trị đo  của từng thời điểm đánh giá so với giá trị trung bình.
	\end{enumerate}
}
\end{bt}

% #1141 | Archive/2024/0. [LVH] BG - 12 New - 4in1-20240625T224301Z-001/0. [LVH] BG - 12 New - 4in1/2. [LVH] TEX SGK 12 NEW-20240418T074833Z-001/2. [LVH] TEX SGK 12 NEW/[LVH] SGK-12-KNTT/data/C3/12-SGK-KNTT-3-10.tex | mã gốc 2D3H2-2
\begin{bt}%[2D3H2-2]
	Thời gian chạy tập luyện cự li $100 \mathrm{~m}$ của hai vận động viên được cho trong bảng sau:
	\begin{center}
		\begin{tabular}{|l|c|c|c|c|}
			\hline Thời gian (giây) & {$[10 ; 10{,}3)$} & {$[10{,}3 ; 10{,}6)$} & {$[10{,}6 ; 10{,}9)$} & {$[10{,}9 ; 11{,}2)$} \\
			\hline Số lần chạy của $A$ & $2$ & $10$ & $5$ & $3$ \\
			\hline Sồ lần chạy của $B$ & $3$ & $7$ & $9$ & $6$ \\
			\hline
		\end{tabular}
	\end{center}	
	Dựa trên độ lệch chuẩn của các mẫu số liệu ghép nhóm, hãy cho biết vận động viên nào có thành tích luyện tập ổn định hơn.
\loigiai{Ta có 	
	\begin{center}
		\begin{tabular}{|l|c|c|c|c|}
			\hline Thời gian (giây) & {$[10 ; 10{,}3)$} & {$[10{,}3 ; 10{,}6)$} & {$[10{,}6 ; 10{,}9)$} & {$[10{,}9 ; 11{,}2)$} \\
			\hline Giá trị đại diện & $10{,}15$ & $10{,}45$ & $10{,}75$ & $11{,}05$ \\
			\hline Số lần chạy của $A$ & $2$ & $10$ & $5$ & $3$ \\
			\hline Sồ lần chạy của $B$ & $3$ & $7$ & $9$ & $6$ \\
			\hline
		\end{tabular}
	\end{center}
	
	\begin{itemize}
		\item Thời gian chạy trung bình củ vận động viên $A$ là $$\overline{x}_A=\dfrac{1}{20}(2\cdot10{,}15+10\cdot10,45+5\cdot10{,}75+3\cdot11{,}05)=10{,}585$$
		\item Độ lệch chuẩn của vận động viên $A$. $$s_A=\sqrt{\dfrac{1}{20}\left(2\cdot10{,}15^2+10\cdot10{,}45^2+5\cdot10{,}75^2+3\cdot11{,}05^2\right)-10{,}585^2}=0{,}25937.$$			
		\item Thời gian chạy trung bình củ vận động viên $B$ là $$\overline{x}_A=\dfrac{1}{25}(3\cdot10{,}15+7\cdot10{,}45+9\cdot10{,}75+6\cdot11{,}05)=10{,}666.$$		
		\item Độ lệch chuẩn của vận động viên $B$. $$s_A=\sqrt{\dfrac{1}{25}\left(3\cdot10{,}15^2+7\cdot10{,}45^2+9\cdot10{,}75^2+6\cdot11{,}05^2\right)-10{,}666^2}=0{,}288.$$		
	\end{itemize}
Vậy độ viên $B$ có thành tích ổn định hơn.
}
\end{bt}

% #1142 | Archive/2024/0. [LVH] BG - 12 New - 4in1-20240625T224301Z-001/0. [LVH] BG - 12 New - 4in1/2. [LVH] TEX SGK 12 NEW-20240418T074833Z-001/2. [LVH] TEX SGK 12 NEW/[LVH] SGK-12-KNTT/data/C3/12-SGK-KNTT-3-10.tex | mã gốc 2D3H2-2
\begin{bt}%[2D3H2-2]
	Có nên dùng phương sai (hoặc độ lệch chuẩn) để so sánh độ phân tán của hai mẫu số liệu ghép nhóm trong mỗi trường hợp sau không? Tại sao?
	\begin{enumerate}[a)]
	\item  Các mẫu số liệu ghép nhóm về điểm thi tốt nghiệp môn Toán của học sinh hai trường trung học phổ thông có chất lượng tương đương.
	\item Các mẫu số liệu ghép nhóm về doanh thu của $100$ cửa hàng bán lẻ và doanh thu của $100$ siêu thị.
	\end{enumerate}
\loigiai{
		\begin{enumerate}[a)]
		\item  Có nên dùng phương sai (hoặc độ lệch chuẩn) để so sánh độ phân tán của học sinh hai trường vì $2$ trường này có chất lượng tương đương nhau
		\item Không nên dùng phương sai (hoặc độ lệch chuẩn) để so sánh độ phân tán về doanh thu của $100$ cửa hàng bán lẻ và $100$ siêu thị vì doanh thu của $100$ cửa hàng bán lẻ và $100$ siêu thị là không gần tương đương nhau.
	\end{enumerate}
}
\end{bt}

% #1143 | Archive/2024/0. [LVH] BG - 12 New - 4in1-20240625T224301Z-001/0. [LVH] BG - 12 New - 4in1/2. [LVH] TEX SGK 12 NEW-20240418T074833Z-001/2. [LVH] TEX SGK 12 NEW/[LVH] SGK-12-KNTT/data/C3/12-SGK-KNTT-3-9.tex | mã gốc 2D3H1-2
\begin{vd}%[2D3H1-2]
	Thống kê thời gian sử dụng mạng xã hội trong ngày của các bạn Tổ 1, Tổ 2 lớp 12A, được kết quả như sau:
	\begin{center}
		\begin{tabular}{|l|c|c|c|c|}
			\hline
			Thời gian sử dụng (phút) & $[0;10 )$ & $[ 10;30)$ &$[30;60 )$ & $[60;90)$ \\
			\hline
			Số học sinh Tổ 1 & $2$ & $4$ & $3$ & $1$ \\
			\hline
			Số học sinh Tổ 1 & $5$ & $1$ & $3$ & $0$ \\
			\hline
		\end{tabular}
	\end{center}
	Tìm khoảng biến thiên cho thời gian sử dụng mạng xã hội của học sinh mỗi tổ và giải thích ý nghĩa.
	\loigiai{
		Gọi $R_1$, $R_2$ tương ứng là khoảng biến thiên của mẫu số liệu ghép nhóm về thời gian sử dụng mạng xã hội trong ngày của các bạn Tổ $1$ và Tổ $2$.
		
		Ta có: $R_1= 90-0=90$ và $R_2= 60-0=60$.
		
		Do $R_1>R_2$ nên ta có thể kết luận rằng thời gian sử dụng mạng xã hội  trong ngày của các bạn Tổ $1$ phân tán hơn thời gian sử dụng  mạng xã hội của các bạn Tổ $2$.
	}
\end{vd}

% #1144 | Archive/2024/0. [LVH] BG - 12 New - 4in1-20240625T224301Z-001/0. [LVH] BG - 12 New - 4in1/2. [LVH] TEX SGK 12 NEW-20240418T074833Z-001/2. [LVH] TEX SGK 12 NEW/[LVH] SGK-12-KNTT/data/C3/12-SGK-KNTT-3-9.tex | mã gốc 2D3H1-3
\begin{vd}%[2D3H1-3]
	Thời gian chờ khám bệnh của các bệnh nhân tại phòng khám X được cho trong bảng sau:
	\begin{center}
		\begin{tabular}{|l|c|c|c|c|}
			\hline
			Thời gian (phút) & $[0;5 )$ & $[5;10)$ &$[10;15 )$ & $[15;20)$ \\
			\hline
			Số bệnh nhân & $3$ & $12$ & $15$ & $8$ \\
			\hline
		\end{tabular}
	\end{center}
	\begin{enumerate}
		\item Tìm khoảng tứ phân vị của mẫu số liệu ghép nhóm này.
		\item Từ một số liệu về thời gian chờ khám bệnh của các bệnh nhân tại phòng khám Y người ta tính được khoảng tứ phần vị bằng $9{,}23$. Hỏi thời gian chờ của bệnh nhân tại phòng khám nào phân tán hơn?
	\end{enumerate}
	\loigiai{
		\begin{enumerate}
			\item Cỡ mẫu là $n = 3+12+15+8=38$. Gọi $x_1,\ldots,x_{38}$ là thời gian chờ khám bệnh của $38$ bệnh nhân này và giả sử rằng dãy số liệu gốc này đã được sắp xếp theo thứ tự tăng dần.
			
			Tứ phân vị thứ nhất của mẫu số liệu gốc là $x_{10}$ nên nhóm chứa tứ phân vị thứ nhất là nhóm $[5;10)$ và ta có:
			\[Q_1 = 5+ \dfrac{\left[ \dfrac{38}{4}-3\right] }{12}\cdot 5 \approx 7{,}71.\]
			Tứ phân vị thứ ba của mẫu số liệu gốc là $x_{29}$ nên nhóm chứa tứ phân vị thứ ba là nhóm $[10;15)$ và ta có:
			\[Q_3 = 10+ \dfrac{\left[ \dfrac{3\cdot 38}{4}-15\right] }{15}\cdot 5 = 14{,}5.\]
			Vậy khoảng tứ phân vị của mẫu số liệu  ghép nhóm là: $\Delta_Q = Q_3 - Q_1 \approx 14{,}5 - 7{,}71= 6{,}79$.
			\item Do $\Delta_Q = 6{,}79 < 9{,}23$ nên thời gian chờ của bệnh nhân tại phòng khám $Y$ phân tán hơn thời gian chờ của bệnh nhân tại phòng khám $X$.
		\end{enumerate}
	}
\end{vd}

% #1145 | Archive/2024/0. [LVH] BG - 12 New - 4in1-20240625T224301Z-001/0. [LVH] BG - 12 New - 4in1/2. [LVH] TEX SGK 12 NEW-20240418T074833Z-001/2. [LVH] TEX SGK 12 NEW/[LVH] SGK-12-KNTT/data/C3/12-SGK-KNTT-3-9.tex | mã gốc 2D3H1-4
\begin{bt}%[2D3H1-4]
	Thống kê số thẻ vàng của mỗi câu lạc bộ trong giải ngoại hạng Anh mùa giải $2021-2022$ cho kết quả sau:	
	\begin{center}
		\begin{tabular}{cccccccccc}
			101&  79&  79&  78&  75&  73&  68&  67&  67& 63 \\
			63&  61&  60&  59&  57&  55&  55&  50&  47& 42.
		\end{tabular}	
		\begin{flushright}
			(Theo  \textit{premierleague.com})
		\end{flushright} 	
	\end{center}
	\begin{enumerate}
		\item Hãy ghép nhóm dãy số liệu trên thành các nhóm có độ dài bằng nhau với nhóm đầu tiên là $[40;50)$.
		\item  Tính khoảng biến thiên, khoảng tứ phân vị của mẫu số liệu gốc và mẫu số liệu ghép nhóm thu được ở câu a. Giá trị nào là giá trị chính xác? Giá trị nào là giá trị xấp xỉ?
	\end{enumerate}
	\loigiai{
		\begin{enumerate}
			\item Ta có bảng dữ liệu sau khi ghép nhóm như sau:
			\begin{center}
				\begin{tabular}{|l|c|c|c|c|c|c|c|}
					\hline
					Số lượng thẻ vàng & $[40;50)$ & $[50;60)$ &$[60;70)$ & $[70;80)$ & $[80;90)$ & $[90;100)$ & $[100;110)$ \\
					\hline
					Tần số & $2$ & $5$ & $7$ & $5$ & $0$& $0$& $1$\\
					\hline					
				\end{tabular}
			\end{center}
			\item Khoảng biến thiên của mẫu số liệu ghép nhóm là 
			\[110-40=70 \text{ (thẻ vàng).}\]
			Với $n=20$, ta có: $\dfrac{n}{4}=5$ mà $2<5<7$. Suy ra nhóm $2$ là nhóm đầu tiên có tần số tích lũy lớn hơn hoặc bằng $5$.\\
			Xét nhóm $2$ là $[50;60)$ có $s=50$, $h=10$, $n_2 = 5$.\\
			Nhóm $1$ là $[40;50)$ có $cf_1=2$ $\Rightarrow Q_1 = 50+\dfrac{5-2}{5}\cdot 10 = 56$.\\
			Ta có: $\dfrac{3n}{4}=15$ mà $14<15<19$ suy ra nhóm $4$ là nhóm có tần số tích lũy lớn hơn hoặc bằng $15$.\\
			Xét nhóm $4$ là $[70;80)$ có $t=70$, $h=10$, $n_4 = 5$.\\
			Nhóm 3 là $[60;70)$ có $cf_3 = 14$ $\Rightarrow Q_3 = 60+\dfrac{15-14}{5}\cdot 10 = 62$.\\
			Vậy $\triangle Q = Q_3-Q_1 = 62-56 = 6$.
			
			
		\end{enumerate}
	}
\end{bt}

% #1146 | Archive/2024/0. [LVH] BG - 12 New - 4in1-20240625T224301Z-001/0. [LVH] BG - 12 New - 4in1/2. [LVH] TEX SGK 12 NEW-20240418T074833Z-001/2. [LVH] TEX SGK 12 NEW/[LVH] SGK-12-KNTT/data/C3/12-SGK-KNTT-3-9.tex | mã gốc 2D3H1-3
\begin{bt}%[2D3H1-3]
	Thu nhập theo tháng (đơn vị: triệu đồng) của người lao động ở hai nhà máy như sau:
	\begin{center}
		\begin{tabular}{|l|c|c|c|c|c|}
			\hline
			Thu nhập & $[5;8 )$ & $[8;11)$ &$[11;14 )$ & $[14;17)$ & $[17;20)$ \\
			\hline
			Số người ở nhà máy A & $20$ & $35$ & $45$ & $35$ & $20$\\
			\hline
			Số người ở nhà máy B & $17$ & $23$ & $30$ & $23$ & $17$\\
			\hline
		\end{tabular}
	\end{center}
	Tính mức thu nhập trung bình của người lao động ở hai nhà máy trên. Dựa vào khoảng tứ phân vị, hãy xác định xem mức thu nhập của người lao động ở nhà máy nào biến động nhiều hơn.
	\loigiai{
		Ta có:
		\begin{center}
			\begin{tabular}{|l|c|c|c|c|c|}
				\hline
				Thu nhập & $[5;8 )$ & $[8;11)$ &$[11;14 )$ & $[14;17)$ & $[17;20)$ \\
				\hline
				Giá trị đại diện & $6,5$ & $9,5$ &$12,5$ & $15,5$ & $18,5$ \\
				\hline
				Số người ở nhà máy A & $20$ & $35$ & $45$ & $35$ & $20$\\
				\hline
				Số người ở nhà máy B & $17$ & $23$ & $30$ & $23$ & $17$\\
				\hline
			\end{tabular}
		\end{center}
		Thu nhập trung bình của người lao động của nhà máy A là
		\[\overline{x}_A = \dfrac{6{,}5\cdot 20 + 9{,}5 \cdot 35 + 12{,}5\cdot45+ 15{,}5\cdot35 +18{,}5\cdot20}{155}= 12{,}5.\]
		Thu nhập trung bình của người lao động của nhà máy B là
		\[\overline{x}_A = \dfrac{6{,}5\cdot 17 + 9{,}5 \cdot 23 + 12{,}5\cdot 30+ 15{,}5\cdot23 +18{,}5\cdot17}{110}= 12{,}5.\]
		Xét nhà máy A, $Q_1 = 10{,}8125$; $Q_3 = 15{,}39$ $\Rightarrow \triangle Q = 4{,}58$.\\
		Xét nhà máy B, $Q_1 = 9{,}85$; $Q_3 = 15{,}63$ $\Rightarrow \triangle Q = 5{,}78$.\\
		Vậy mức thu nhập của người lao động ở nhà máy B có biến động nhiều hơn
	}
\end{bt}

% #1147 | Archive/2024/0. [LVH] BG - 12 New - 4in1-20240625T224301Z-001/0. [LVH] BG - 12 New - 4in1/2. [LVH] TEX SGK 12 NEW-20240418T074833Z-001/2. [LVH] TEX SGK 12 NEW/[LVH] SGK-12-KNTT/data/C3/12-SGK-KNTT-3-9.tex | mã gốc 2D3H1-4
\begin{bt}%[2D3H1-4]
	Bảng sau đây cho biết chiều cao của các học sinh lớp 12A và 12B.
	\begin{center}
		\begin{tabular}{|l|c|c|c|c|c|c|}
			\hline
			Chiều cao (cm) & $[145;150 )$ & $[150;155)$ &$[155;160 )$ & $[160;165)$ & $[165;170)$ & $[170;175)$ \\
			\hline
			Số học sinh của lớp 12A & $1$ & $0$ & $15$ & $12$ & $10$ & $5$\\
			\hline
			Số học sinh của lớp 12B & $0$ & $0$ & $17$ & $10$ & $9$ & $6$\\
			\hline
		\end{tabular}
	\end{center}
	\begin{enumerate}
		\item  Tính khoảng biến thiên, khoảng tứ phần vị cho các mẫu số liệu ghép nhóm của học sinh lớp 12A, 12B.
		\item Để so sánh độ phân tán về chiều cao của học sinh hai lớp này ta nên dùng khoảng biến thiên hay khoảng tứ phân vị? Vì sao?		
	\end{enumerate}
	\loigiai{
	\begin{enumerate}
		\item Ta có
		\begin{center}
			\begin{tabular}{|l|c|c|c|c|c|c|}
				\hline
				Chiều cao (cm) & $[145;150 )$ & $[150;155)$ &$[155;160 )$ & $[160;165)$ & $[165;170)$ & $[170;175)$ \\
				\hline
				Số học sinh của lớp 12A & $1$ & $0$ & $15$ & $12$ & $10$ & $5$\\
				\hline
				Số học sinh của lớp 12B & $0$ & $0$ & $17$ & $10$ & $9$ & $6$\\
				\hline
			\end{tabular}
		\end{center}
		Khoảng biến thiên là $175-145 = 30$ (cm).\\
		Xét lớp 12A,\\
		\[Q_1 = 155 + \dfrac{\dfrac{43}{4}-1}{15}\cdot 5 = 158{,}25.\]
		\[Q_3 = 165 + \dfrac{\dfrac{43\cdot 3}{4}-28}{10}\cdot 5 = 167{,}125.\]
		\[\triangle Q = Q_3 -Q_1 = 8{,}875.\]
		Xét lớp 12B,\\
		\[Q_1 = 155 + \dfrac{\dfrac{42}{4}-0}{17}\cdot 5 = 158{,}5\]
		\[Q_3 = 165 + \dfrac{\dfrac{42\cdot 3}{4}-27}{9}\cdot 5 = 167{,}5\]
		\[\triangle Q = Q_3 -Q_1 = 9{,}4.\]
		\item Để so sánh độ phân tán về chiều cao của học sinh hai lớp này ta nên dùng khoảng tứ phân vị, vì khoảng biến thiên của $2$ lớp này là bằng nhau.
	\end{enumerate}
	}
\end{bt}

% #1901 | Archive/2024/TaiLieuGiangDay/3. K12 - 4in1 - Thongke/data/2. Bai1-Dang2-KhoangTuPhanVi.tex | mã gốc 2D3N1-3
\begin{bt}[Mức độ 1]%[2D3N1-3]
	Chiều cao của $500$ học sinh của một trường trung học cơ sở được thống kê bởi bảng dưới đây
	\begin{center}
		{\setlength\extrarowheight{2pt}
			\begin{tabular}{|>{\centering\arraybackslash}m{3cm}|*{5}{>{\centering\arraybackslash}m{2cm}|}}
				\hline 
				Chiều cao (cm) & $[150; 154)$ & $ [154; 158)$ & $ [158; 162)$ & $ [162; 166)$ & $ [166; 170)$   \\
				\hline
				Số học sinh &$25$ & $50$ & $200$ & $175$ & $50$ \\
				\hline 
		\end{tabular}}
	\end{center}
	Tính khoảng tứ phân vị. Kết quả tìm được cho biết điều gì về chiều cao của $500$ học sinh này?
	\loigiai
	{Ta có cỡ mẫu $n= 500$ và bảng tần số tích lũy như sau
		\begin{center}
			{\setlength\extrarowheight{2pt}
				\begin{tabular}{|>{\centering\arraybackslash}m{3cm}|*{5}{>{\centering\arraybackslash}m{2cm}|}}
					\hline 
					Nhóm & $[150; 154)$ & $ [154; 158)$ & $ [158; 162)$ & $ [162; 166)$ & $ [166; 170)$   \\
					\hline
					Tần số &$25$ & $50$ & $200$ & $175$ & $50$ \\
					\hline 
					Tần số tích lũy & $25$ & $ 75 $ & $275$ & $450$ & $500$ \\
					\hline
			\end{tabular}}
		\end{center}
		Ta có $\dfrac{n}{4}= 125$, suy ra nhóm chứa $Q_1$ là $[158; 162)$. Do đó $Q_1 = 158 + \dfrac{125- 75}{200} \cdot 4 = 159$.\\
		Lại có $\dfrac{3n}{4}= 375$, suy ra nhóm chứa $Q_3$ là $[162; 166)$. Do đó $Q_3= 162 +\dfrac{375-275}{175} \cdot 4 =\dfrac{1150}{7} = \approx 164{,}3$.\\
		Vậy khoảng tứ phân vị của mẫu số liệu là $\Delta _Q= Q_3 - Q_1 = \dfrac{37}{7} \approx 5{,}3$.
		Dựa vào các tứ phân vị trên ta thấy rằng trong mẫu số liệu có khoảng $50\%$ học sinh có chiều cao từ $159$ cm đến $164{,}3$ cm.
	}
\end{bt}

% #1902 | Archive/2024/TaiLieuGiangDay/3. K12 - 4in1 - Thongke/data/2. Bai1-Dang2-KhoangTuPhanVi.tex | mã gốc 2D3N1-3
\begin{bt}[Mức độ 1]%[2D3N1-3]
	Một trang trại phân $1 \, 000$ quả trứng thành $5$ loại, tuỳ theo khối lượng (đã được làm tròn) của chúng được thống kê bởi bảng dưới đây:
	\begin{center}
		{\setlength\extrarowheight{2pt}
			\begin{tabular}{|>{\centering\arraybackslash}m{3.5cm}|*{5}{>{\centering\arraybackslash}m{1.9cm}|}}
				\hline 
				Khối lượng (gam) & $[30; 36)$ & $ [36; 42)$ & $ [42; 48)$ & $ [48; 54)$ & $ [54; 60)$   \\
				\hline
				Số trứng & $45$ & $190$ & $500$ & $250$ & $15$ \\
				\hline 
		\end{tabular}}
	\end{center}
	Hãy tính khoảng tứ phân vị của mẫu số liệu trên.
	\loigiai{
		Cỡ mẫu là $n= 1\,000$. Ta có bảng tần số tích lũy của mẫu số liệu như sau
		\begin{center}
			{\setlength\extrarowheight{2pt}
				\begin{tabular}{|>{\centering\arraybackslash}m{3.5cm}|*{5}{>{\centering\arraybackslash}m{1.9cm}|}}
					\hline 
					Nhóm & $[30; 36)$ & $ [36; 42)$ & $ [42; 48)$ & $ [48; 54)$ & $ [54; 60)$   \\
					\hline
					Tần số & $45$ & $190$ & $500$ & $250$ & $15$ \\
					\hline 
					Tần số tích lũy & $45$ & $235$ & $735$ & $985$ & $1\, 000$ \\
					\hline
			\end{tabular}}
		\end{center}
		Ta có $\dfrac{n}{4}= 250$ nên nhóm chứa $Q_1$ là nhóm $[42; 48)$. Suy ra 
		$Q_1= 42 + \dfrac{250- 235}{500} \cdot 6=42{,}18$.\\
		Lại có $\dfrac{3n}{4}= 750$ nên nhóm chứa $Q_3$ là nhóm $[48; 54)$. Do đó $Q_3 =48 +\dfrac{750- 735}{250} \cdot 6 = 48{,}36$.\\
		Vậy khoảng tứ phân vị của mẫu số liệu là $\Delta_Q = Q_3 - Q_1= 6{,}18$.}
\end{bt}

% #1903 | Archive/2024/TaiLieuGiangDay/3. K12 - 4in1 - Thongke/data/2. Bai1-Dang2-KhoangTuPhanVi.tex | mã gốc 2D3N1-3
\begin{bt}[Mức độ 1]%[BG12-New-4in1, Nguyễn Trung Kiên]%[2D3N1-3]
	Dưới đây là kết quả điều tra thời gian hoàn thành bài tập ở nhà môn Toán của 30 học sinh lớp 10 một trường phổ thông A
	\begin{center}
		\textit{Thời gian hoàn thành bài tập môn Toán ở nhà (đơn vị: phút)}
		{\setlength\extrarowheight{2pt}
		\begin{tabular}
			{|*{10}{>{\centering\arraybackslash}m{1cm}|}}
			\hline 
			$42$ & $45$ & $47$ & $47$ & $53$ & $54$ & $58$ & $58$ & $58$ & $59$ \\
			\hline 
			$61$ & $63$ & $64$ & $64$ & $67$ & $68$ & $68$ & $68$ & $70$ & $73$ \\
			\hline 
			$75$ & $75$ & $77$ & $77$ & $78$ & $78$ & $82$ & $82$ & $82$ & $87$ \\
			\hline
		\end{tabular}}
	\end{center}	
	\begin{enumerate}
		\item Lập mẫu số liệu ghép nhóm với các nhóm ghép có độ dài bằng $10$ và nhóm đầu tiên là $[40;50)$. Xác định giá trị đại diện và tần số tích lũy của các nhóm.
		\item Tìm số trung bình cộng, khoảng biến thiên, các tứ phân vị và khoảng tứ phân vị của mẫu số liệu ghép nhóm vừa lập.
	\end{enumerate}
	\loigiai{
		\begin{enumerate}
			\item Ta có $n=30$, suy ra bảng mẫu số liệu ghép nhóm như sau:
			\begin{center}
				{\setlength\extrarowheight{2pt}
				\begin{tabular}{|*{4}{>{\centering\arraybackslash}m{3.5cm}|}}
					\hline 
					Nhóm & Giá trị đại diện & Tần số & Tần số tích lũy\\
					\hline$[40 ; 50)$ & $45$ & $4$ & $4$ \\
					\hline$[50 ; 60)$ & $55$ & $6$ & $10$ \\
					\hline$[60 ; 70)$ & $65$ & $8$ & $18$ \\
					\hline$[70 ; 80)$ & $75$ & $8$ & $26$ \\
					\hline$[80 ; 90)$ & $85$ & $4$ & $30$ \\
					\hline
				\end{tabular}}	
			\end{center}
			\item Ta có
			\begin{itemize}
				\item Giá trị trung bình $\bar{x}=\dfrac{45 \cdot 4+55 \cdot 6+65 \cdot 8+75 \cdot 8+85 \cdot 4}{30} = \dfrac{197}{3}\approx 65{,}67$.
				\item Khoảng biến thiên $R=90-40=50$.
				\item Ta có $n=30$, $\dfrac{n}{4}=7{,}5$, suy ra $Q_1$ thuộc nhóm $[50;60)$. Từ đó ta có $$Q_1=50+\dfrac{7{,}5 -4}{6} \cdot 10 = \dfrac{335}{6}.$$
				Lại có $\dfrac{n}{2}=15$, suy ra $Q_2$ thuộc nhóm $[60;70)$. Từ đó ta có $$Q_2=60+\dfrac{15-10}{8} \cdot 10= \dfrac{265}{4}.$$
				Ta có $\dfrac{3n}{4}=22{,}5$, suy ra $Q_3$ thuộc nhóm $[70;80)$. Vì vậy 
				$$Q_3=70+\dfrac{22{,}5 -18}{8} \cdot 10 = \dfrac{605}{8}.$$
				\item Vậy khoảng tứ phân vị là $\Delta_Q = Q_3-Q_1 \approx \dfrac{605}{8} -\dfrac{335}{6} \approx 19{,}79$.
			\end{itemize}
		\end{enumerate}
	}
\end{bt}

% #1904 | Archive/2024/TaiLieuGiangDay/3. K12 - 4in1 - Thongke/data/2. Bai1-Dang2-KhoangTuPhanVi.tex | mã gốc 2D3N1-3
\begin{bt}[Mức độ 1]%[2D3N1-3]
	Thời gian trung bình hằng ngày mà một số nhân viên đi từ nhà đến công ty được thống kê trong bảng sau:
	\begin{center}
		{\setlength\extrarowheight{2pt}
		\begin{tabular}{|>{\centering\arraybackslash}m{3.5cm}|*{6}{>{\centering\arraybackslash}m{1.6cm}|}}
			\hline 
			Thời gian (phút)  & {$[10 ; 20)$} & {$[20 ; 30)$} & {$[30 ; 40)$} & {$[40 ; 50)$} & {$[50 ; 60)$} & {$[60 ; 70)$} \\
			\hline 
			Số người (tần số)  & $8$ & $8$ & $10$ & $6$ & $7$ & $1$ \\
			\hline
		\end{tabular}}
	\end{center}
	Tìm khoảng biến thiên và khoảng tứ phân vị của thời gian di chuyển đến công ty của các nhân viên (làm tròn kết quả đến hàng phần mười).	
	\loigiai{
		Ta có bảng tần số tích luỹ như sau:
		\begin{center}
			{\setlength\extrarowheight{2pt}
				\begin{tabular}{|>{\centering\arraybackslash}m{3.5cm}|*{6}{>{\centering\arraybackslash}m{1.6cm}|}}
					\hline 
					Nhóm & {$[10 ; 20)$} & {$[20 ; 30)$} & {$[30 ; 40)$} & {$[40 ; 50)$} & {$[50 ; 60)$} & {$[60 ; 70)$} \\
					\hline 
					Tần số & $8$ & $8$ & $10$ & $6$ & $7$ & $1$ \\
					\hline
					Tần số tích luỹ & $8$ & $16$ & $26$ & $32$ & $39$ & $40$ \\
					\hline
			\end{tabular}}
		\end{center}
		\begin{itemize}
			\item Khoảng biến thiên $R=70-10=60$.
			\item Ta có kích thước mẫu là $n=40$.\\
			Vì $\dfrac{n}{4}=10$ nên nhóm chứa $Q_1$ là nhóm $[20;30)$, suy ra
			$$Q_1=20+\dfrac{10 -8}{8} \cdot 10= \dfrac{45}{2}.$$
			Mặt khác $\dfrac{3n}{4}=30$ nên nhóm chứa $Q_3$ là nhóm $[40;50)$, suy ra
			$$Q_3=40+\dfrac{30 -26}{6} \cdot 10 =\dfrac{140}{3}.$$
			Vậy khoảng tứ phân vị là $\Delta_Q =\dfrac{140}{3} -\dfrac{45}{2} = \dfrac{145}{6} \approx 24{,}2$.
		\end{itemize}
	}
\end{bt}

% #1905 | Archive/2024/TaiLieuGiangDay/3. K12 - 4in1 - Thongke/data/2. Bai1-Dang2-KhoangTuPhanVi.tex | mã gốc 2D3H1-3
\begin{bt}[Mức độ 2]%[2D3H1-3]
	Bảng số liệu sau trình bày dữ liệu về tốc độ của $100$ xe ô tô lưu thông trên một đoạn đường cao tốc vào giờ cao điểm, được trích xuất từ camera của cơ quan cảnh sát giao thông.
	\begin{center}
		{\setlength\extrarowheight{2pt}
			\begin{tabular}{|>{\centering\arraybackslash}m{3.5cm}|*{5}{>{\centering\arraybackslash}m{1.8cm}|}}
				\hline 
				Tốc độ (km/h) & $[60 ; 70)$& $[70 ; 80)$ & $[80 ; 90)$ & $[90 ; 100)$ & $[100 ; 110)$ \\
				\hline 
				Số xe & $10$ & $20$ & $20$ & $35$ & $15$ \\
				\hline
		\end{tabular}}
	\end{center}
	Hãy tìm khoảng tứ phân vị của mẫu số liệu (làm tròn kết quả đến hàng đơn vị). Nêu ý nghĩa của kết quả tìm được.
	\loigiai{
		\begin{itemize}
			\item Ta có bảng tần số tích luỹ như sau:
			\begin{center}
				{\setlength\extrarowheight{2pt}
					\begin{tabular}{|>{\centering\arraybackslash}m{3.5cm}|*{5}{>{\centering\arraybackslash}m{1.8cm}|}}
						\hline 
						Nhóm & $[60 ; 70)$& $[70 ; 80)$ & $[80 ; 90)$ & $[90 ; 100)$ & $[100 ; 110)$ \\
						\hline 
						Tần số & $10$ & $20$ & $20$ & $35$ & $15$ \\
						\hline
						Tần số tích luỹ & $10$ & $30$ & $50$ & $85$ & $100$ \\
						\hline
				\end{tabular}}
			\end{center}
			Ta có kích thước mẫu là $n=100$.\\
			Vì $\dfrac{n}{4}=25$ nên nhóm chứa $Q_1$ là $[70;80)$, suy ra
			$$Q_1=70+\dfrac{25 -10}{20} \cdot 10 = \dfrac{155}{2}.$$
			Mặt khác $\dfrac{3n}{4}=75$ nên nhóm chứa $Q_3$ là $[90;100)$, suy ra
			$$Q_3=90+\dfrac{75 -50}{35} \cdot 10 =\dfrac{680}{7}.$$
			Vậy khoảng tứ phân vị là $\Delta_Q = \dfrac{680}{7} -\dfrac{155}{2} = \dfrac{275}{14} \approx 97-78 \approx 20$. 
			\item Ý nghĩa: Khoảng $50\%$ lượng xe lưu thông đi trên đường cao tốc giờ cao điểm chạy tốc độ từ $77{,}5$ km/h đến $97{,}1$ km/h. 
		\end{itemize}		
	}
\end{bt}

% #1906 | Archive/2024/TaiLieuGiangDay/3. K12 - 4in1 - Thongke/data/2. Bai1-Dang2-KhoangTuPhanVi.tex | mã gốc 2D3H1-3
\begin{bt}[Mức độ 2]%[2D3H1-3]
	Một ngân hàng thống kê ở bảng dưới số tiền (nghìn đồng) mà khách hàng rút từ một máy ATM (máy rút tiền tự động) trong một buổi sáng.
	\begin{center}
		{\setlength\extrarowheight{2pt}
			\begin{tabular}{|>{\centering\arraybackslash}m{3cm}|*{6}{>{\centering\arraybackslash}m{2.1cm}|}}
				\hline 
				Số tiền rút & $[0;500)$ & $[500;1\,000)$ & $[1\,000;1\,500)$ & $[1\,500;2\,000)$ & $[2\,000;2\,500)$ & $[2\,500;3\,000)$\\
				\hline 
				Tần số & $11$ & $16$ & $12$ & $15$ & $10$ & $16$ \\
				\hline
		\end{tabular}}
	\end{center}
	Tìm khoảng tứ phân vị của số tiền rút (làm tròn kết quả đến hàng đơn vị). Nêu ý nghĩa của các kết quả tìm được.
	\loigiai{
		Ta có bảng tần số tích luỹ như sau:
		\begin{center}
			{\setlength\extrarowheight{2pt}
				\begin{tabular}{|>{\centering\arraybackslash}m{3cm}|*{6}{>{\centering\arraybackslash}m{2.1cm}|}}
					\hline 
					Nhóm & $[0;500)$ & $[500;1\,000)$ & $[1\,000;1\,500)$ & $[1\,500;2\,000)$ & $[2\,000;2\,500)$ & $[2\,500;3\,000)$\\
					\hline 
					Tần số & $11$ & $16$ & $12$ & $15$ & $10$ & $16$ \\
					\hline
					Tần số tích lũy & $11$ & $27$ & $39$ & $54$ & $64$ & $80$ \\
					\hline
			\end{tabular}}
		\end{center}
		Ta có kích thước mẫu là $n=80$.\\
		Vì $\dfrac{n}{4}=20$ nên nhóm chứa $Q_1$ là nhóm $[500;1 \,000)$, suy ra
		$$Q_1=500+\dfrac{20 -11}{16} \cdot 500 \approx 781.$$
		Mặt khác $\dfrac{3n}{4}=60$ nên nhóm chứa $Q_3$ là nhóm $[2\,000;2\,500)$, suy ra
		$$Q_3=2\,000+\dfrac{60 -54}{10} \cdot 500=2\,300.$$
		Vậy khoảng tứ phân vị là $\Delta_Q \approx 2 \,300-781 \approx 1\, 519$.\\
		Ý nghĩa: Có khoảng $50\%$ số khách hàng rút tiền dao động từ $781$ nghìn đồng đến $2$ triệu $300$ nghìn đồng.}
\end{bt}

% #1907 | Archive/2024/TaiLieuGiangDay/3. K12 - 4in1 - Thongke/data/2. Bai1-Dang2-KhoangTuPhanVi.tex | mã gốc 2D3H1-3
\begin{bt}[Mức độ 2]%[2D3H1-3]
	Hàm lượng protein (trong $100$ g) của một số loại thực phẩm được cho trong bảng sau:
	\begin{center}
		{\setlength\extrarowheight{2pt}
			\begin{tabular}{|>{\centering\arraybackslash}m{4.2cm}|*{6}{>{\centering\arraybackslash}m{1.6cm}|}}
				\hline 
				Hàm lượng protein(g) & $[8; 10)$ & $ [10; 12)$ & $ [12; 14)$ & $ [14; 16)$ & $ [16; 18]$ & $[18;20]$   \\
				\hline
				Tần số &$4$ & $12$ & $16$ & $14$ & $2$ & $2$ \\
				\hline 
		\end{tabular}}
	\end{center}
	Tìm khoảng tứ phân vị của mẫu số liệu (làm tròn kết quả đến hàng phần trăm). Nêu ý nghĩa của các kết quả tìm được.
	\loigiai{
		Cỡ mẫu $n= 4+ 12+ 16+ 14+ 2 + 2= 50$. Ta có bảng tần số tích lũy của mẫu số liệu như sau:
		\begin{center}
			{\setlength\extrarowheight{2pt}
				\begin{tabular}{|>{\centering\arraybackslash}m{4.2cm}|*{6}{>{\centering\arraybackslash}m{1.6cm}|}}
					\hline 
					Nhóm & $[8; 10)$ & $ [10; 12)$ & $ [12; 14)$ & $ [14; 16)$ & $ [16; 18]$ & $[18;20]$   \\
					\hline
					Tần số &$4$ & $12$ & $16$ & $14$ & $2$ & $2$ \\
					\hline 
					Tần số tích lũy &$4$ & $16$ & $32$ & $46$ & $48$ & $50$ \\
					\hline
			\end{tabular}}
		\end{center}
		Vì $\dfrac{n}{4}=12{,}5$ nên nhóm chứa $Q_1$ là nhóm $[10; 12)$. Suy ra $Q_1= 10 + \dfrac{12{,}5- 4}{12} \cdot 2= \dfrac{137}{12}\approx 11{,} 42$.\\
		Mặt khác $\dfrac{3n}{4}= 37{,}5$ nên nhóm chứa $Q_3$ là nhóm $[14; 16)$. Suy ra $Q_3 =14 +\dfrac{37{,}5- 32 }{14} \cdot 2 =\dfrac{207}{14} \approx 14{,}79$.\\
		Suy ra khoảng tứ phân vị của mẫu số liệu là $\Delta_Q = Q_3 - Q_1 = \dfrac{283}{84}\approx  3{,}37$.\\
		Ý nghĩa: Ta có khoảng $25\%$ loại thực phẩmcó hàm lượng protein không vượt quá $11{,} 42$ gam, $25\%$ loại thực phẩm có hàm lượng protein đạt ít nhất $14{,} 79$ gam và có khoảng $50\%$ loại thực phẩm còn lại có hàm lượng protein từ $11{,} 42$ gam đến $14{,}79$ gam.}
\end{bt}

% #1920 | Archive/2024/TaiLieuGiangDay/3. K12 - 4in1 - Thongke/data/4. Bai2-Dang1-PhuongSai,Dolechchuan.tex | mã gốc 2D3H2-2
\begin{vd}[Mức độ 3]%[Ngô Quang Anh]%[2D3H2-2]
	Mai và Ngọc cùng sử dụng vòng đeo tay thông minh để ghi lại số bước chân hai bạn đi mỗi ngày trong một tháng. Kết quả được ghi lại ở bảng sau
	\begin{center}
		\begin{tabular}{|c|c|c|c|c|c|}
			\hline Số bước (đơn vị: nghìn) &{$[3;5)$} &{$[5;7)$} &{$[7;9)$} &{$[9;11)$} &{$[11;13)$} \\
			\hline Mai & $ 6 $ & $ 7 $ & $ 6 $ & $ 6 $ & $ 5 $ \\
			\hline Ngọc & $ 2 $ & $ 5 $ & $ 13 $ & $ 8 $ & $ 2 $ \\
			\hline
		\end{tabular}
	\end{center}
	\begin{enumerate}
		\item Hãy tính số trung bình và độ lệch chuẩn của mẫu số liệu ghép nhóm trên.
		\item Nếu so sánh theo độ lệch chuẩn thì bạn nào có số lượng bước chân đi mỗi ngày đều đặn hơn?
	\end{enumerate}
	\loigiai{
		\begin{enumerate}
			\item Ta có bảng sau
			\begin{center}
				\begin{tabular}{|c|c|c|c|c|c|}
					\hline Số bước (đơn vị: nghìn) &{$[3;5)$} &{$[5;7)$} &{$[7;9)$} &{$[9;11)$} &{$[11;13)$} \\
					\hline Số bước đại diện & $ 4 $ & $ 6 $ & $ 8 $ & $ 10 $ & $ 12 $ \\
					\hline Mai & $ 6 $ & $ 7 $ & $ 6 $ & $ 6 $ & $ 5 $ \\
					\hline Ngọc & $ 2 $ & $ 5 $ & $ 13 $ & $ 8 $ & $ 2 $ \\
					\hline
				\end{tabular}
			\end{center}
			\begin{itemize}
				\item Xét mẫu số liệu của Mai.\\
				Cỡ mẫu là $ n_1=6+7+6+6+5=30 $.\\
				Số trung bình của mẫu số liệu ghép nhóm là\\
				$$ \overline{x}_1=\dfrac{6\cdot 4+7\cdot6+6\cdot8+6\cdot10+5\cdot12}{30}=7{,}8. $$
				Phương sai của mẫu số liệu ghép nhóm là
				$$ S_1^2=\dfrac{1}{30}\left(6\cdot 4^2+7\cdot6^2+6\cdot8^2+6\cdot10^2+5\cdot12^2 \right) -7{,}8^2=7{,}56. $$
				Độ lệch chuẩn của mẫu số liệu ghép nhóm là $ S_1=\sqrt{7{,}56}.$
				
				\item Xét mẫu số liệu của Ngọc.\\
				Cỡ mẫu là $ n_2=2+5+13+8+2=30 $.\\
				Số trung bình của mẫu số liệu ghép nhóm là
				$$ \overline{x}_2=\dfrac{2\cdot 4+5\cdot6+13\cdot8+8\cdot10+2\cdot12}{30}=8{,}2. $$
				Phương sai của mẫu số liệu ghép nhóm là
				$$ S_2^2=\dfrac{1}{30}\left(2\cdot 4^2+5\cdot6^2+13\cdot8^2+8\cdot10^2+2\cdot12^2 \right) -8{,}2^2=\dfrac{287}{75}. $$
				Độ lệch chuẩn của mẫu số liệu ghép nhóm là $ S_2=\sqrt{\dfrac{287}{75}}. $
			\end{itemize}
			\item Do $ S_1>S_2 $ nên nếu so sánh theo độ lệch chuẩn thì Ngọc có số lượng bước chân đi mỗi ngày đều đặn hơn Mai.
		\end{enumerate}
	}
\end{vd}

% #1921 | Archive/2024/TaiLieuGiangDay/3. K12 - 4in1 - Thongke/data/4. Bai2-Dang1-PhuongSai,Dolechchuan.tex | mã gốc 2D3H2-2
\begin{vd}[Mức độ 3]%[Ngô Quang Anh]%[2D3H2-2]
	Giá đóng cửa của một cổ phiếu là giá của cổ phiếu đó cuối một phiên giao dịch. Bảng sau thống kê giá đóng cửa (đơn vị: nghìn đồng) của hai mã cổ phiếu $A$ và $B$ trong $50$ ngày giao dịch liên tiếp.
	\begin{center}
		\begin{tabular}{|c|c|c|c|c|c|}
			\hline Giá đóng cửa &{$[120; 122)$} &{$[122; 124)$} &{$[124; 126)$} &{$[126; 128)$} &{$[128; 130)$} \\
			\hline Cổ phiếu $ A $ & $ 8 $ & $ 9 $ & $ 12 $ & $ 10 $ & $ 11 $ \\
			\hline Cổ phiếu $ B $ & $ 16 $ & $ 4 $ & $ 3 $ & $ 6 $ & $ 21 $ \\
			\hline
		\end{tabular}
	\end{center}
	Hãy so sánh độ rủi ro của cổ phiếu $A$ và cổ phiếu $B$.
	\loigiai{
		Ta có bảng thống kê giá đóng cửa theo giá trị đại diện
		\begin{center}
			\begin{tabular}{|c|c|c|c|c|c|}
				\hline Giá đóng cửa &{$121$} &{$123$} &{$125$} &{$127$} &{$129$} \\
				\hline Cổ phiếu $ A $ & $ 8 $ & $ 9 $ & $ 12 $ & $ 10 $ & $ 11 $ \\
				\hline Cổ phiếu $ B $ & $ 16 $ & $ 4 $ & $ 3 $ & $ 6 $ & $ 21 $ \\
				\hline
			\end{tabular}
		\end{center}
		\begin{itemize}
			\item Xét mẫu số liệu của cổ phiếu A\\
			Số trung bình của mẫu số liệu ghép nhóm là
			$$ \overline{x}_1=\dfrac{8\cdot121+9\cdot123+12\cdot125+10\cdot127+11\cdot129}{50}=125{,}28. $$
			Phương sai của mẫu số liệu ghép nhóm là
			$$ S_1^2=\dfrac{1}{50}\left(8\cdot121^2+9\cdot123^2+12\cdot125^2+10\cdot127^2+11\cdot129^2 \right)-125{,}28^2=7{,}5216.  $$
			Độ lệch chuẩn của mẫu số liệu ghép nhóm là $ S_1=\sqrt{7{,}5216}. $
			\item Xét mẫu số liệu của cổ phiếu B\\
			Số trung bình của mẫu số liệu ghép nhóm là
			$$ \overline{x}_2=\dfrac{16\cdot121+4\cdot123+3\cdot125+6\cdot127+21\cdot129}{50}=125{,}48. $$
			Phương sai của mẫu số liệu ghép nhóm là
			$$ S_2^2=\dfrac{1}{50}\left(16\cdot121^2+4\cdot123^2+3\cdot125^2+6\cdot127^2+21\cdot129^2 \right)-125{,}48^2=12{,}4096.  $$
			Độ lệch chuẩn của mẫu số liệu ghép nhóm là $ S_2=\sqrt{12{,}4096}. $
		\end{itemize}
		Vì $S_1<S_2$ nên cổ phiếu $A$ có độ rủi ro thấp hơn cổ phiếu $B$.
	}
\end{vd}

% #1922 | Archive/2024/TaiLieuGiangDay/3. K12 - 4in1 - Thongke/data/4. Bai2-Dang1-PhuongSai,Dolechchuan.tex | mã gốc 2D3H2-2
\begin{vd}[Mức độ 3]%[Ngô Quang Anh]%[2D3H2-2]
	Thống kê lợi nhuận hàng tháng (đơn vị: triệu đồng) trong $20$ tháng của hai nhà đầu tư được cho như sau:
	\begin{center}
		\begin{tabular}{|l|c|c|c|c|c|}
			\hline Lợi nhuận & {$[10; 20)$} & {$[20; 30)$} & {$[30; 40)$} & {$[40; 50)$} & {$[50 ; 60)$} \\
			\hline Số tháng & $2$ & $4$ & $8$ & $4$ & $2$ \\
			\hline
		\end{tabular}
	\end{center}	
	\begin{center}
		Bảng 3.2. Lợi nhuận theo tháng của nhà đầu tư nhỏ
	\end{center}
	\begin{center}
		\begin{tabular}{|l|c|c|c|c|c|}
			\hline Lợi nhuận & {$[510; 520)$} & {$[520; 530)$} & {$[530; 540)$} & {$[540; 550)$} & {$[550; 560)$} \\
			\hline Số tháng & $4$ & $3$ & $6$ & $3$ & $4$ \\
			\hline
		\end{tabular}\\
		Bảng 3.3. Lợi nhuận theo tháng của nhà đầu tư lớn
	\end{center}	
	Tính độ lệch chuẩn của hai mẫu số liệu ghép nhóm trên. Có nên dựa vào độ lệch chuẩn để so sánh độ rủi ro của hai nhà đầu tư này không?
	\loigiai{
		Chọn điểm đại diện cho các nhóm số liệu ta tính được các số đặc trưng như sau\\
		Lợi nhuận trung bình một tháng của các nhà đầu tư tương ứng là
		$$
		\bar{x}_A=\dfrac{1}{20}(2 \cdot 15+\ldots+2 \cdot 55)=35 \text { (triệu đồng); } \bar{x}_B=\dfrac{1}{20}(4 \cdot 515+\ldots+4 \cdot 555)=535 \text { (triệu đồng). }
		$$
		Độ lệch chuẩn của lợi nhuận hàng tháng của hai nhà đầu tư tương ứng là
		$$s_A=\sqrt{\dfrac 1{20}\left(2\cdot 15^2+\ldots+2\cdot 55^2\right)-(35)^2}\approx 10{,}95;$$
		$$s_B=\sqrt{\dfrac 1{20}\left(4\cdot 515^2+\ldots+4\cdot 555^2\right)-(535)^2}\approx 13{,}78.$$
		Độ lệch chuẩn cho lợi nhuận hàng tháng của nhà đầu tư lớn cao hơn của nhà đầu tư nhỏ. Lợi nhuận trung bình của hai nhà đầu tư khác nhau rất nhiều, do đó ta không nên dùng độ lệch chuẩn để so sánh mức độ rủi ro của hai nhà đầu tư này.
	}
\end{vd}

% #1925 | Archive/2024/TaiLieuGiangDay/3. K12 - 4in1 - Thongke/data/4. Bai2-Dang1-PhuongSai,Dolechchuan.tex | mã gốc 2D3H2-2
\begin{bt}[Mức độ 2]%[Dự án bài giảng 12, Nguyễn Tấn Linh]%[2D3H2-2]
Số lượng sách đọc của sinh viên trong một tháng được thống kê ở bảng sau:
\begin{center}
\begin{tabular}{|c|c|c|c|c|}
\hline
\textbf{Số sách (quyển)} & $[0; 2)$ & $[2; 4)$ & $[4; 6)$ & $[6; 8)$ \\
\hline
\textbf{Số sinh viên} & $8$ & $15$ & $22$ & $5$ \\
\hline
\end{tabular}
\end{center}
Hãy tính phương sai của mẫu số liệu ghép nhóm này.
\loigiai
{
Ta có bảng thống kê số sách đọc của sinh viên theo giá trị đại diện
\begin{center}
\begin{tabular}{|c|c|c|c|c|}
\hline
\textbf{Số sách đại diện (quyển)} & 1 & 3 & 5 & 7 \\
\hline
\textbf{Số sinh viên} & 8 & 15 & 22 & 5 \\
\hline
\end{tabular}
\end{center}
Cỡ mẫu $n=8+15+22+5=50$.\\
Số trung bình của mẫu số liệu ghép nhóm là:
\[ \overline{x}=\dfrac{8\cdot 1+15\cdot 3+22\cdot 5+5\cdot 7}{50}=3{,}96. \]
Phương sai của mẫu số liệu ghép nhóm là:
\[ S^2=\dfrac{1}{50}\left(8\cdot 1^2+15\cdot 3^2+22\cdot 5^2+5\cdot 7^2 \right)-3{,}96^2= 3{,}0784. \]
}
\end{bt}

% #1926 | Archive/2024/TaiLieuGiangDay/3. K12 - 4in1 - Thongke/data/4. Bai2-Dang1-PhuongSai,Dolechchuan.tex | mã gốc 2D3H2-2
\begin{bt}[Mức độ 2]%[Dự án bài giảng 12, Nguyễn Tấn Linh]%[2D3H2-2]
Chiều cao của một số học sinh trong một lớp được thống kê ở bảng sau:
\begin{center}
\begin{tabular}{|c|c|c|c|c|c|}
\hline
\textbf{Chiều cao (cm)} & $[140; 150)$ & $[150; 160)$ & $[160; 170)$ & $[170; 180)$ & $[180; 190)$ \\
\hline
\textbf{Số học sinh} & $5$ & $10$ & $15$ & $8$ & $2$ \\
\hline
\end{tabular}
\end{center}
Hãy tính phương sai của mẫu số liệu ghép nhóm này.
\loigiai
{
Ta có bảng thống kê chiều cao của các học sinh theo giá trị đại diện:
\begin{center}
\begin{tabular}{|c|c|c|c|c|c|}
\hline
\textbf{Chiều cao đại diện (cm)} & 145 & 155 & 165 & 175 & 185 \\
\hline
\textbf{Số học sinh} & 5 & 10 & 15 & 8 & 2 \\
\hline
\end{tabular}
\end{center}
Cỡ mẫu $n=5+10+15+8+2=40$.\\
Số trung bình của mẫu số liệu ghép nhóm là:
\[ \overline{x}=\dfrac{5\cdot 145+10\cdot155+15\cdot165+8\cdot175+2\cdot185}{40}=163. \]
Phương sai của mẫu số liệu ghép nhóm là:
\[ S^2=\dfrac{1}{40}\left(5\cdot145^2+10\cdot 155^2+15\cdot165^2+8\cdot175^2+2\cdot185^2 \right)-163^2=111. \]
}
\end{bt}

% #1927 | Archive/2024/TaiLieuGiangDay/3. K12 - 4in1 - Thongke/data/4. Bai2-Dang1-PhuongSai,Dolechchuan.tex | mã gốc 2D3H2-2
\begin{bt}[Mức độ 2]%[Dự án bài giảng 12, Nguyễn Tấn Linh]%[2D3H2-2]
Nhiệt độ trung bình hàng ngày trong một tuần được thống kê ở bảng sau:
\begin{center}
\begin{tabular}{|c|c|c|c|c|c|c|c|}
\hline
\textbf{Nhiệt độ ($^\circ$C)} & $[20; 22)$ & $[22; 24)$ & $[24; 26)$ & $[26; 28)$ & $[28; 30)$ & $[30; 32)$ & $[32; 34)$ \\
\hline
\textbf{Số ngày} & $1$ & $1$ & $2$ & $1$ & $0$ & $1$ & $1$ \\
\hline
\end{tabular}
\end{center}
Hãy tính phương sai của mẫu số liệu ghép nhóm này.
\loigiai
{
Ta có bảng thống kê nhiệt độ trung bình hàng ngày theo giá trị đại diện:
\begin{center}
\begin{tabular}{|c|c|c|c|c|c|c|c|}
\hline
\textbf{Nhiệt độ đại diện ($^\circ$C)} & $21$ & $23$ & $25$ & $27$ & $29$ & $31$ & $33$ \\
\hline
\textbf{Số ngày} & $1$ & $1$ & $2$ & $1$ & $0$ & $1$ & $1$ \\
\hline
\end{tabular}
\end{center}
Cỡ mẫu $n=1+1+2+1+0+1+1=7$.\\
Số trung bình của mẫu số liệu ghép nhóm là:
\[ \overline{x}=\dfrac{1\cdot 21+1\cdot 23+2\cdot 25+1\cdot 27+0\cdot 29+1\cdot 31+1\cdot 33}{7}=26{,}42. \]
Phương sai của mẫu số liệu ghép nhóm là:
\[ S^2=\dfrac{1}{7}\left(1\cdot 21^2+1\cdot 23^2+2\cdot 25^2+1\cdot 27^2+0\cdot 29^2+1\cdot 31^2+1\cdot 33^2 \right)-26{,}42^2\approx 15{,}67. \]
}
\end{bt}

% #1928 | Archive/2024/TaiLieuGiangDay/3. K12 - 4in1 - Thongke/data/4. Bai2-Dang1-PhuongSai,Dolechchuan.tex | mã gốc 2D3H2-2
\begin{bt}[Mức độ 2]%[Dự án bài giảng 12, Nguyễn Tấn Linh]%[2D3H2-2]
Thời gian chạy 100m của các vận động viên trong một giải đấu được thống kê ở bảng sau:
\begin{center}
\begin{tabular}{|c|c|c|c|c|}
\hline
\textbf{Thời gian (s)} & $[10; 11)$ & $[11; 12)$ & $[12; 13)$ & $[13; 14)$ \\
\hline
\textbf{Số vận động viên} & $3$ & $7$ & $12$ & $8$ \\
\hline
\end{tabular}
\end{center}
Hãy tính phương sai của mẫu số liệu ghép nhóm này.
\loigiai
{
Ta có bảng thống kê tốc độ của các vận động viên theo giá trị đại diện:
\begin{center}
\begin{tabular}{|c|c|c|c|c|}
\hline
\textbf{Thời gian đại diện (s)} & $10{,}5$ & $11{,}5$ & $12{,}5$ & $13{,}5$ \\
\hline
\textbf{Số vận động viên} & $3$ & $7$ & $12$ & $8$ \\
\hline
\end{tabular}
\end{center}
Cỡ mẫu $n=3+7+12+8=30$.\\
Số trung bình của mẫu số liệu ghép nhóm là
\[ \overline{x}=\dfrac{3\cdot 10{,}5+7\cdot 11{,}5+12\cdot 12{,}5+8\cdot 13{,}5}{30}=12{,}33. \]
Phương sai của mẫu số liệu ghép nhóm là:
\[ S^2=\dfrac{1}{30}\left(3\cdot 10{,}5^2+7\cdot 11{,}5^2+12\cdot 12{,}5^2+8\cdot 13{,}5^2 \right)-12{,}33^2\approx 0{,}87. \]
}
\end{bt}

% #1929 | Archive/2024/TaiLieuGiangDay/3. K12 - 4in1 - Thongke/data/4. Bai2-Dang1-PhuongSai,Dolechchuan.tex | mã gốc 2D3H2-2
\begin{bt}[Mức độ 3]%[Dự án bài giảng 12, Nguyễn Tấn Linh]%[2D3H2-2]
Một loại cây cảnh được trồng tại hai khu vực khác nhau $C$ và $D$ trong công viên. Người ta thống kê chiều cao của một số cây cảnh 3 năm tuổi ở bảng sau
\begin{center}
\begin{tabular}{|c|c|c|c|c|c|}
\hline Chiều cao (cm)  & {$[50; 55)$} & {$[55; 60)$} & {$[60; 65)$} & {$[65; 70)$} & {$[70; 75)$} \\
\hline Số cây trồng ở khu vực $ C $  & $ 15 $ & $ 25 $ & $ 30 $ & $ 20 $ & $ 10 $ \\
\hline Số cây trồng ở khu vực $ D $  & $ 10 $ & $ 20 $ & $ 25 $ & $ 35 $ & $ 10 $ \\
\hline
\end{tabular}
\end{center}
\begin{enumerate}
\item Hãy so sánh chiều cao trung bình của cây cảnh trồng tại khu vực $C$ và khu vực $D$.
\item Nếu so sánh theo phương sai thì cây trồng tại khu vực nào có chiều cao đồng đều hơn?
\end{enumerate}
\loigiai{
\begin{enumerate}
\item Ta có bảng sau
\begin{center}
\begin{tabular}{|c|c|c|c|c|c|}
\hline Chiều cao (cm)  & {$[50; 55)$} & {$[55; 60)$} & {$[60; 65)$} & {$[65; 70)$} & {$[70; 75)$} \\
\hline Chiều cao đại diện  & $ 52.5 $ & $ 57.5 $ & $ 62.5 $ & $ 67.5 $ & $ 72.5 $ \\
\hline Số cây trồng ở khu vực $ C $  & $ 15 $ & $ 25 $ & $ 30 $ & $ 20 $ & $ 10 $ \\
\hline Số cây trồng ở khu vực $ D $  & $ 10 $ & $ 20 $ & $ 25 $ & $ 35 $ & $ 10 $ \\
\hline
\end{tabular}
\end{center}
Cỡ mẫu $ n=100 $.
\begin{itemize}
\item Chiều cao trung bình của cây cảnh trồng ở khu vực $C$
$$ \overline{x}_C=\dfrac{15\cdot52.5+25\cdot57.5+30\cdot62.5+20\cdot67.5+10\cdot72.5}{100}=61.75. $$
\item Chiều cao trung bình của cây cảnh trồng ở khu vực $D$
$$ \overline{x}_D=\dfrac{10\cdot52.5+20\cdot57.5+25\cdot62.5+35\cdot67.5+10\cdot72.5}{100}=63.25.$$
\end{itemize}
Ta có $\overline{x}_D>\overline{x}_C$ nên chiều cao trung bình cây cảnh ở khu vực $ D $ lớn hơn ở khu vực $ C $.
\item \begin{itemize}
\item Phương sai của mẫu số liệu cây cảnh trồng ở khu vực $ C $
$$ S_C^2=\dfrac{1}{100}\left(15\cdot52.5^2+25\cdot57.5^2+30\cdot62.5^2+20\cdot67.5^2+10\cdot72.5^2 \right)- 61.75^2=43.1875. $$
\item Phương sai của mẫu số liệu cây cảnh trồng ở khu vực $ D $
$$ S_D^2=\dfrac{1}{100}\left(10\cdot52.5^2+20\cdot57.5^2+25\cdot62.5^2+35\cdot67.5^2+10\cdot72.5^2 \right)- 63.25^2=38.6875. $$
\end{itemize}
Ta có $ S^2_D<S^2_C $ nên nếu so sánh theo độ lệch chuẩn thì cây trồng tại khu vực $ D $ đồng đều hơn ở khu vực $ C $.
\end{enumerate}
}
\end{bt}

% #1930 | Archive/2024/TaiLieuGiangDay/3. K12 - 4in1 - Thongke/data/4. Bai2-Dang1-PhuongSai,Dolechchuan.tex | mã gốc 2D3H2-2
\begin{bt}[Mức độ 3]%[Ngô Quang Anh]%[2D3H2-2]
	Bảng sau thống kê lại tổng số giờ nắng trong tháng $6$ của các năm từ $2002$ đến $2021$ tại hai trạm quan trắc đặt ở Nha Trang và Quy Nhơn
	\begin{center}
		\begin{tabular}{|c|c|c|c|c|c|c|}
			\hline
			Số giờ nắng &$[130;160)$  & $[160;190)$ & $[190;220)$ & $[220;250)$ & $[250;280)$ & $[280;310)$\\
			\hline
			Số năm ở Nha Trang	& $1$ & $1$ & $1$ & $8$ & $7$ & $2$\\
			\hline
			Số năm ở Quy Nhơn & $0$ & $1$ & $2$ & $4$ & $10$ & $3$\\
			\hline
		\end{tabular}
	\end{center}
	Hỏi số giờ nắng trong tháng $6$ của địa phương nào đồng đều hơn?
	\loigiai 
	{
		Ta có bảng thống kê tổng số giờ nắng trong tháng 6 ở hai địa phương theo giá trị đại diện.
		\begin{center}
			\begin{tabular}{|c|c|c|c|c|c|c|}
				\hline
				Số giờ nắng &$[130;160)$  & $[160;190)$ & $[190;220)$ & $[220;250)$ & $[250;280)$ & $[280;310)$\\
				\hline
				Giá trị đại diện &$145$  & $175$ & $205$ & $235$ & $265$ & $295$\\
				\hline
				Số năm ở Nha Trang	& $1$ & $1$ & $1$ & $8$ & $7$ & $2$ \\
				\hline
				Số năm ở Quy Nhơn & $0$ & $1$ & $2$ & $4$ & $10$ & $3$ \\
				\hline
			\end{tabular}
		\end{center}
		\begin{itemize}
			\item Xét mẫu số liệu ở Nha Trang.\\
			Cỡ mẫu là $n_X=20$.\\
			Số trung bình của mẫu số liệu ghép nhóm là 
			$$\overline{x}_X=\dfrac{1}{20}\left(1\cdot 145+1\cdot 175+1\cdot 205+8\cdot 235+7\cdot 265+2\cdot 295\right)=242{,}5.$$
			Phương sai của mẫu số liệu ghép nhóm là
			$$S^2_X=\dfrac{1}{20}\left(1\cdot 145^2+1\cdot 175^2+1\cdot 205^2+8\cdot 235^2+7\cdot 265^2+2\cdot 295^2\right)-242{,}5^2=1248{,}75.$$
			Độ lệch chuẩn của mẫu số liệu là 
			$$S_X=\sqrt{1248{,}75}\approx 33{,}338.$$
			\item Xét mẫu số liệu ở Quy Nhơn.\\
			Cỡ mẫu là $n_Y=20$.\\
			Số trung bình của mẫu số liệu ghép nhóm là 
			$$\overline{x}_Y=\dfrac{1}{20}\left(0\cdot 145+1\cdot 175+2\cdot 205+4\cdot 235+10\cdot 265+3\cdot 295\right)=253.$$
			Phương sai của mẫu số liệu ghép nhóm là
			$$S^2_Y=\dfrac{1}{20}\left(0\cdot 145^2+1\cdot 175^2+2\cdot 205^2+4\cdot 235^2+107\cdot 265^2+3\cdot 295^2\right)-253^2=936.$$
			Độ lệch chuẩn của mẫu số liệu là 
			$$S_Y=\sqrt{936}\approx 30{,}594.$$
		\end{itemize}
		Vì $S_Y<S_X$ nên số giờ nắng trong tháng 6 của Quy Nhơn đồng đều hơn.
	}
\end{bt}

% #1931 | Archive/2024/TaiLieuGiangDay/3. K12 - 4in1 - Thongke/data/4. Bai2-Dang1-PhuongSai,Dolechchuan.tex | mã gốc 2D3V2-2
\begin{bt}[Mức độ 3]%[Ngô Quang Anh]%[2D3V2-2]
	Biểu đồ sau mô tả kết quả điều tra về điểm trung bình năm học của học sinh hai trường $A$ và $B$.
	\begin{center}
		\begin{tikzpicture}[line join=round, line cap=round,>=stealth,xscale=1,yscale=1]
			\def\a{1}
			\def\xmax{12}
			\def\ymax{8}
			\tikzset{label style/.style={font=\footnotesize}}
			\draw[->] (0,0)--(\xmax,0) node[below] {\text{Điểm trung bình}};
			\draw[->] (0,0)--(0,\ymax) node[below left] {\text{Số học sinh}};
			\draw (0,0) node [below left] {$0$};
			\draw[thin,gray!60]
			(0,1)--(\xmax,1)
			(0,2)--(\xmax,2)
			(0,3)--(\xmax,3)
			(0,4)--(\xmax,4)
			(0,5)--(\xmax,5)
			(0,6)--(\xmax,6)
			(0,7)--(\xmax,7)
			;
			\foreach \i in {1,2,3,4,5,6,7}{
				\draw (0,\i) node[left]{$\i$};
			}
			\foreach \i/\j in {1*\a/[5;6),3*\a/{[6;7)},5*\a/{[7;8)},7*\a/{[8;9)},9*\a/{[9;10)}}{
				\draw (\i,-.5) node[below]{\small $\j$};
			}
			\fill[gray!50]
			(0,0) rectangle (\a,4)
			(2*\a,0) rectangle (3*\a,5)
			(4*\a,0) rectangle (5*\a,3)
			(6*\a,0) rectangle (7*\a,4)
			(8*\a,0) rectangle (9*\a,2);
			\fill[blue!30!]
			(\a,0) rectangle (2*\a,2)
			(3*\a,0) rectangle (4*\a,5)
			(5*\a,0) rectangle (6*\a,4)
			(7*\a,0) rectangle (8*\a,3)
			(9*\a,0) rectangle (10*\a,1);
			\begin{scope}
				\draw 
				(0,4)--(\a,4) 
				(\a,2)--(2*\a,2)
				(2*\a,5)--(3*\a,5)
				(3*\a,5)--(4*\a,5)
				(4*\a,3)--(5*\a,3)
				(5*\a,4)--(6*\a,4)
				(6*\a,4)--(7*\a,4)
				(7*\a,3)--(8*\a,3)
				(8*\a,2)--(9*\a,2)
				(9*\a,1)--(10*\a,1)
				(0,0)--(0,4)
				(\a,0)--(\a,4)
				(2*\a,0)--(2*\a,5)
				(3*\a,0)--(3*\a,5)
				(4*\a,0)--(4*\a,5)
				(5*\a,0)--(5*\a,4)
				(6*\a,0)--(6*\a,4)
				(7*\a,0)--(7*\a,4)
				(8*\a,0)--(8*\a,3)
				(9*\a,0)--(9*\a,2)
				(10*\a,0)--(10*\a,1);
				\draw (6,\ymax+1) node[above]{\textbf{Điểm trung bình năm học của học sinh hai trường A và B}};
			\end{scope}
			\foreach \col/\t [count=\i from 0] in {gray!50/học sinh trường A, blue!30!/học sinh trường B}\path (10,6-\i) node[right]{\t} node[left,fill=\col]{};
		\end{tikzpicture}
	\end{center}
	Hỏi học sinh trường nào có điểm trung bình đồng đều hơn?
	\loigiai 
	{
		Ta có bảng thống kê điểm trung bình của học sinh ở hai trường theo giá trị đại diện.
		\begin{center}
			\begin{tabular}{|c|c|c|c|c|c|}
				\hline
				Điểm trung bình & $[5;6)$ &$[6;7)$  & $[7;8)$ & $[8;9)$ & $[9;10)$  \\
				\hline
				Giá trị đại diện & $5{,}5$ & $6{,}5$ & $7{,}5$ & $8{,}5$ & $9{,}5$  \\
				\hline
				Học sinh trường $A$	& $4$ & $5$ & $3$ & $4$ & $2$ \\
				\hline
				Học sinh trường $B$	& $2$ & $5$ & $4$ & $3$ & $1$ \\
				\hline
			\end{tabular}
		\end{center}
		\begin{itemize}
			\item Xét mẫu số liệu học sinh trường $A$.\\
			Cỡ mẫu là $n_A=18$.\\
			Số trung bình của mẫu số liệu ghép nhóm là 
			$$\overline{x}_A=\dfrac{1}{18}\left(4\cdot 5{,}5+5\cdot 6{,}5+3\cdot 7{,}5+4\cdot 8{,}5+2\cdot 9{,}5\right)=\dfrac{65}{9}.$$
			Phương sai của mẫu số liệu ghép nhóm là
			$$S^2_A=\dfrac{1}{18}\left(4\cdot 5{,}5^2+5\cdot 6{,}5^2+3\cdot 7{,}5^2+4\cdot 8{,}5^2+2\cdot 9{,}5^2\right)-\left(\dfrac{65}{9}\right)^2=1{,}756.$$
			Độ lệch chuẩn của mẫu số liệu là 
			$$S_A=\sqrt{1{,}756}\approx 1{,}325.$$
			\item Xét mẫu số liệu học sinh trường $B$.\\
			Cỡ mẫu là $n_B=15$.\\
			Số trung bình của mẫu số liệu ghép nhóm là 
			$$\overline{x}_B=\dfrac{1}{15}\left(2\cdot 5{,}5+5\cdot 6{,}5+4\cdot 7{,}5+3\cdot 8{,}5+1\cdot 9{,}5\right)=\dfrac{217}{30}.$$
			Phương sai của mẫu số liệu ghép nhóm là
			$$S^2_B=\dfrac{1}{15}\left(2\cdot 5{,}5^2+5\cdot 6{,}5^2+4\cdot 7{,}5^2+3\cdot 8{,}5^2+1\cdot 9{,}5^2\right)-\left(\dfrac{217}{30}\right)^2=1{,}262.$$
			Độ lệch chuẩn của mẫu số liệu là 
			$$S_B=\sqrt{1{,}262}\approx 1{,}123.$$
		\end{itemize}
		Vì $S_B<S_A$ nên học sinh trường $B$ có điểm trung bình đồng đều hơn.
	}
\end{bt}

% #1932 | Archive/2024/TaiLieuGiangDay/3. K12 - 4in1 - Thongke/data/4. Bai2-Dang1-PhuongSai,Dolechchuan.tex | mã gốc 2D3V2-2
\begin{bt}[Mức độ 3]%[Ngô Quang Anh]%[2D3V2-2]
	Kết quả 40 lần nhảy xa của hai vận động viên nam Dũng và Huy được lần lượt thống kê trong các bảng sau (đơn vị: mét)
	\begin{center}
		\begin{tabular}{|c|c|}
			\hline Nhóm & Tần số \\
			\hline$[6,22 ; 6,46)$ & 3 \\
			{$[6,46 ; 6,70)$} & 7 \\
			{$[6,70 ; 6,94)$} & 5 \\
			{$[6,94 ; 7,18)$} & 20 \\
			{$[7,18 ; 7,42)$} & 5 \\
			\hline & $n=40$ \\
			\hline
		\end{tabular}
		\begin{tabular}{|c|c|}
			\hline Nhóm & Tần số \\
			\hline$[6,22 ; 6,46)$ & 2 \\
			{$[6,46 ; 6,70)$} & 5 \\
			{$[6,70 ; 6,94)$} & 8 \\
			{$[6,94 ; 7,18)$} & 19 \\
			{$[7,18 ; 7,42)$} & 6 \\
			\hline & $n=40$ \\
			\hline
		\end{tabular}
	\end{center}
	\begin{enumerate}
		\item  Tính độ lệch chuẩn của mẫu số liệu ghép nhóm biểu diễn kết quả 40 lần nhảy xa của vận động viên Dũng cho bởi Bảng thứ nhất (làm tròn kết quả đến hàng phần trăm).
		\item Tính độ lệch chuẩn của mẫu số liệu ghép nhóm biểu diễn kết quả 40 lần nhảy xa của vận động viên Huy cho bởi Bảng thứ hai (làm tròn kết quả đến hàng phẩn trăm).
		\item Trong hai vận động viên đó, kết quả nhảy xa của vận động viên nào đồng đều hơn?
	\end{enumerate}
	\loigiai{
		Ta có các bảng thống kê sau:
		\begin{center}
			\begin{tabular}{|c|c|c|}
				\hline Nhóm & \begin{tabular}{c} 
					Giá trị \\
					đại diện
				\end{tabular} & Tần số \\
				\hline$[6,22 ; 6,46)$ & 6,34 & 3 \\
				{$[6,46 ; 6,70)$} & 6,58 & 7 \\
				{$[6,70 ; 6,94)$} & 6,82 & 5 \\
				{$[6,94 ; 7,18)$} & 7,06 & 20 \\
				{$[7,18 ; 7,42)$} & 7,30 & 5 \\
				\hline & & $n=40$ \\
				\hline
			\end{tabular}
			\begin{tabular}{|c|c|c|}
				\hline Nhóm & \begin{tabular}{c} 
					Giá trị \\
					đại diện
				\end{tabular} & Tần số \\
				\hline$[6,22 ; 6,46)$ & 6,34 & 2 \\
				{$[6,46 ; 6,70)$} & 6,58 & 5 \\
				{$[6,70 ; 6,94)$} & 6,82 & 8 \\
				{$[6,94 ; 7,18)$} & 7,06 & 19 \\
				{$[7,18 ; 7,42)$} & 7,30 & 6 \\
				\hline & & $n=40$ \\
				\hline
			\end{tabular}
		\end{center}
		\begin{enumerate}
			\item Số trung bình cộng của mẫu số liệu ghép nhóm biểu diễn kết quå 40 lần nhảy xa của vận động viên Dũng cho bởi Bảng thứ nhất là:
			$$
			\bar{x}_D=\frac{3 \cdot 6,34+7 \cdot 6,58+5 \cdot 6,82+20 \cdot 7,06+5 \cdot 7,30}{40}=\frac{276,88}{40} \approx 6,92(\text{m}) .
			$$
			
			Vậy phương sai của mẫu số liệu ghép nhóm biểu diễn kết quả 40 lần nhảy xa của vận động viên Dũng cho bởi Bảng thứ nhất (làm tròn kết quả đến hàng phần trăm) là:
			$$
			\begin{aligned}
				s_D^2=\frac{1}{40}\left[3 .(6,34-6,92)^2\right. & +7 \cdot(6,58-6,92)^2+5 .(6,82-6,92)^2 \\
				& \left.+20 \cdot(7,06-6,92)^2+5 \cdot(7,30-6,92)^2\right]=\frac{2,9824}{40} \approx 0,07 .
			\end{aligned}
			$$
			
			Độ lệch chuẩn của mẫu số liệu ghép nhóm trên là: $s_D \approx \sqrt{0,07} \approx 0,26(\text{m})$.
			\item Số trung bình cộng của mẫu số liệu ghép nhóm biểu diễn kết quả 40 lần nhảy xa của vận động viên Huy cho bởi Bảng thứ hai là:
			$$
			\bar{x}_H=\frac{2 \cdot 6,34+5 \cdot 6,58+8 \cdot 6,82+19 \cdot 7,06+6 \cdot 7,30}{40}=\frac{278,08}{40} \approx 6,95(\text{m}) .
			$$
			$$
			\begin{aligned}
				s_H^2 & =\frac{1}{40}\left[2 \cdot(6,34-6,95)^2+5 \cdot(6,58-6,95)^2\right. \\
				& +8 \cdot(6,82-6,95)^2+19 \cdot(7,06-6,95)^2 \\
				& \left.+6 \cdot(7,30-6,95)^2\right]=\frac{2,5288}{40} \approx 0,06 .
			\end{aligned}
			$$
			
			Độ lệch chuẩn của mẫu số liệu ghép nhóm trên là:
			$$
			s_H \approx \sqrt{0,06} \approx 0,24(\text{m}) .
			$$
			\item Do $s_H \approx 0,24<s_D \approx 0,26$ nên kết quả nhảy xa của vận động viên Huy đồng đều hơn kết quả nhảy xa của vận động viên Dũng.
			
		\end{enumerate}
	}
\end{bt}

% #1955 | Archive/2024/TaiLieuGiangDay/3. K12 - 4in1 - Thongke/data/1. Bai1-Dang1-Khoangbienthien.tex | mã gốc 2D3N1-2
\begin{vd}[Sách Cánh diều, mức độ 1]%[2D3N1-2]
	\immini{Bảng 1 biểu diễn mẫu số liệu ghép nhóm về chiều cao (đơn vị: centimet) của 36 học sinh nam lớp $12$ ở một trường trung học phổ thông. Tìm khoảng biến thiên của mẫu số liệu ghép nhóm đó.}{\begin{tabular}{|c|c|}
			\hline
			Nhóm & Tần số\\
			\hline
			$[160; 163)$ & 6 \\
			$[163; 166)$ & 11 \\
			$[166; 169)$ & 9 \\
			$[169; 172)$ & 7 \\
			$[172; 175)$ & 3 \\
			\hline
			& $n=36$ \\
			\hline
			\multicolumn{2}{c}{Bảng 1} \\
	\end{tabular}}
	\loigiai{
		Trong mẫu số liệu ghép nhóm đó, ta có: đầu mút trái của nhóm $1$ là $a_{1} = 160$, đầu mút phải của nhóm 5 là $a_6 = 175$. Vậy khoảng biến thiên của mẫu số liệu ghép nhóm đó là $R=a_6 - a_1=175-160=15~(\mathrm{cm})$.
	}			
\end{vd}

% #1956 | Archive/2024/TaiLieuGiangDay/3. K12 - 4in1 - Thongke/data/1. Bai1-Dang1-Khoangbienthien.tex | mã gốc 2D3N1-2
\begin{vd}[Sách Chân trời sáng tạo, mức độ 1]%[2D3N1-2]
	Bạn Trang thống kê lại chiều cao (đơn vị: cm) của các bạn học sinh nữ lớp 12C và lớp 12D
	ở bảng sau.
	\begin{center}
		\begin{tabular}{|c|c|c|c|c|c|c|}
			\hline
			\diagbox[width=8em]{Lớp}{Chiều cao} & $[155;160)$ & $[160;165)$ & $[165;170)$ & $[170;175)$ & $[175;180)$ & $[180;185)$ \\
			\hline
			12C & $2$ & $7$ & $12$ & $3$ & $0$ & $1$ \\
			\hline
			12D & $5$ & $9$ & $8$ & $2$ & $1$ & $0$ \\
			\hline
		\end{tabular}
	\end{center}	
	Nếu so sánh theo khoảng biến thiên thì chiều cao của học sinh lớp nào có độ phân tán lớn hơn?
	\loigiai{\hfill\\
		Khoảng biến thiên của chiều cao học sinh nữ lớp 12C là $R_C=185-155=30$.\\
		Khoảng biến thiên của chiều cao học sinh nữ lớp 12D là $R_D=180-155=25$.\\
		Nhận thấy $R_C>R_D$ nên chiều cao của học sinh lớp 12C có độ phân tán lớn hơn.}
\end{vd}

% #1959 | Archive/2024/TaiLieuGiangDay/3. K12 - 4in1 - Thongke/data/1. Bai1-Dang1-Khoangbienthien.tex | mã gốc 2D3N1-2
\begin{bt}[Mức độ 1]%[2D3N1-2]
	Cô Hà thống kê lại đường kính thân gỗ của một số cây xoan đào 6 năm tuổi được trồng ở một lâm trường ở bảng sau.
	\begin{center}
		\begin{tabular}{|c|c|c|c|c|c|}
			\hline
			Đường kính (cm) & $[40;45)$ & $[45;50)$ & $[50;55)$ & $[55;60)$ & $[60;65)$ \\
			\hline
			Tần số & $5$ & $20$ & $18$ & $7$ & $3$ \\
			\hline
		\end{tabular}
	\end{center}
	Hãy tìm khoảng biến thiên của mẫu số liệu ghép nhóm trên.
	\loigiai{Khoảng biến thiên của mẫu số liệu ghép nhóm trên là $R=65-40=25$ (cm).}
\end{bt}

% #1960 | Archive/2024/TaiLieuGiangDay/3. K12 - 4in1 - Thongke/data/1. Bai1-Dang1-Khoangbienthien.tex | mã gốc 2D3N1-2
\begin{bt}[Mức độ 1]%[2D3N1-2]
	Bảng sau thống kê tổng lượng mưa (đơn vị: mm) đo được vào tháng 7 từ năm 2002 đến 2021 tại một trạm quan trắc đặt ở Cà Mau.
	\begin{center}
		\begin{tabular}{|c|c|c|c|c|c|c|c|c|c|}
			\hline
			$341{,}4$ & $187{,}1$ & $242{,}2$ & $522{,}9$ & $251{,}4$ & $432{,}2$ & $200{,}7$ & $388{,}6$ & $258{,}4$ & $288{,}5$ \\
			\hline
			$298{,}1$ & $413{,}5$ & $413{,}5$ & $332$ & $421$ & $475$ & $400$ & $305$ & $520$ & $147$ \\
			\hline
		\end{tabular}\\
		(Nguồn: Tổng cục Thống kê)
	\end{center}
	Hãy tìm khoảng biến thiên của mẫu số liệu trên.
	\loigiai{Xếp lại mẫu số liệu theo thứ tự không giảm, ta được
		\begin{center}
			\begin{tabular}{|c|c|c|c|c|c|c|c|c|c|}
				\hline
				$147$ & $187{,}1$ & $200{,}7$ & $242{,}2$ & $251{,}4$ & $258{,}4$ & $288{,}5$ & $298{,}1$ & $305$ & $332$ \\
				\hline
				$341{,}4$ & $388{,}6$ & $400$ & $413{,}5$ & $413{,}5$ & $421$ & $432{,}2$ & $475$ & $520$ & $522{,}9$ \\
				\hline
			\end{tabular}
		\end{center}
		Khoảng biến thiên của mẫu số liệu trên là $R=522{,}9-147=375{,}9$.
	}
\end{bt}

% #1961 | Archive/2024/TaiLieuGiangDay/3. K12 - 4in1 - Thongke/data/1. Bai1-Dang1-Khoangbienthien.tex | mã gốc 2D3N1-2
\begin{bt}[Mức độ 1]%[2D3N1-2]
	Kết quả đo chiều cao của $100$ cây keo 3 năm tuổi tại một nông trường được cho ở bảng sau
	\begin{center}
		\begin{tabular}{|c|c|c|c|c|c|}
			\hline
			Chiều cao (m) & $[8{,}4;8{,}6)$ & $[8{,}6;8{,}8)$ & $[8{,}8;9{,}0)$ & $[9{,}0;9{,}2)$ & $[9{,}2;9{,}4)$ \\
			\hline
			Số cây & $5$ & $12$ & $25$ & $44$ & $14$ \\
			\hline
		\end{tabular}
	\end{center}
	Hãy tìm khoảng biến thiên của mẫu số liệu ghép nhóm trên.
	%	\end{enumerate}
\loigiai{Khoảng biến thiên của mẫu số liệu là $R = 9{,}4 - 8{,}4 = 1$.
}
\end{bt}

% #1962 | Archive/2024/TaiLieuGiangDay/3. K12 - 4in1 - Thongke/data/1. Bai1-Dang1-Khoangbienthien.tex | mã gốc 2D3N1-2
\begin{bt}[Mức độ 1]%[2D3N1-2]
	{\it Bảng 3.3} thống kê thành tích nhảy xa của một số học sinh lớp $12$. Tìm khoảng biến thiên thành tích nhảy xa của số học sinh này.
	\begin{center}
		{\it Bảng 3.3. Thành tích nhảy xa của một số học sinh lớp $12$}\\
		\begin{tabular}{|p{4cm}|p{2cm}|p{2cm}|p{2cm}|p{2cm}|p{2cm}|}
			\hline Thành tích $(\mathrm{cm})$&\centering\arraybackslash {$[150; 180)$} &\centering\arraybackslash {$[180; 210)$} &\centering\arraybackslash {$[210; 240)$} &\centering\arraybackslash {$[240; 270)$} &\centering\arraybackslash {$[270; 300)$} \\
			\hline Số học sinh &\centering\arraybackslash $3$ &\centering\arraybackslash $5$ &\centering\arraybackslash $28$ &\centering\arraybackslash $14$ &\centering\arraybackslash $8$ \\
			\hline
		\end{tabular}
	\end{center}
	\loigiai{
		Đầu mút phải của nhóm ghép cuối cùng là $300$, đầu mút trái của nhóm ghép đầu tiên là $150$. Vậy khoảng biến thiên của mẫu số liệu là $R=300-150=150$.
	}
\end{bt}

% #1963 | Archive/2024/TaiLieuGiangDay/3. K12 - 4in1 - Thongke/data/1. Bai1-Dang1-Khoangbienthien.tex | mã gốc 2D3N1-2
\begin{bt}[Mức độ 1]%[2D3N1-2]
	\textit{Bảng 3.10} trình bày dữ liệu về tốc độ của $100$ xe ô tô lưu thông trên một đoạn đường cao tốc vào giờ cao điểm, được trích xuất từ camera của cơ quan cảnh sát giao thông. Hãy tìm khoảng biến thiên của mẫu số liệu (làm tròn kết quả đến hàng đơn vị).
	\begin{center}
		\textit{Bảng 3.10}\\
		\begin{tabular}{|l|c|c|c|c|c|}
			\hline  \textbf{Tốc độ (km/h)} & $[60 ; 70)$& $[70 ; 80)$ & $[80 ; 90)$ & $[90 ; 100)$ & $[100 ; 110)$ \\
			\hline \textbf{Số xe} & $10$ & $20$ & $20$ & $35$ & $15$ \\
			\hline
		\end{tabular}
	\end{center} 
	\loigiai{Khoảng biến thiên $R = 110 - 50 = 50$.
	}
\end{bt}

% #1964 | Archive/2024/TaiLieuGiangDay/3. K12 - 4in1 - Thongke/data/1. Bai1-Dang1-Khoangbienthien.tex | mã gốc 2D3H1-2
\begin{bt}[Mức độ 2]%[2D3H1-2]
{\it Bảng 3.2} biểu thị kết quả điều tra thời gian sử dụng Internet hằng ngày của một số người.
\begin{center}
	{\it Bảng 3.2. Thời gian sử dụng Internet hằng ngày}
	\begin{tabular}{|p{4cm}|p{2cm}|p{2cm}|p{2cm}|p{2cm}|p{2cm}|}
		\hline Thời gian (phút) &\centering\arraybackslash {$[30;60)$} &\centering\arraybackslash {$[60;90)$} &\centering\arraybackslash {$[90;120)$} &\centering\arraybackslash {$[120;150)$} &\centering\arraybackslash {$[150;180)$} \\
		\hline Số người &\centering\arraybackslash $2$ &\centering\arraybackslash $4$ &\centering\arraybackslash $10$ &\centering\arraybackslash $5$ &\centering\arraybackslash $3$ \\
		\hline
	\end{tabular}
\end{center}
Tìm khoảng biến thiên của mẫu số liệu đã cho. Kết quả cho biết điều gì?
\loigiai{Đầu mút phải của nhóm ghép cuối cùng là $180$, đầu mút trái của nhóm ghép đầu tiên là $30$. Vậy khoảng biến thiên của mẫu số liệu là $R = 180 - 30 = 150$.\\
	Kết quả này cho biết thời gian sử dụng Internet hằng ngày của các thành viên thuộc nhóm người được điều tra chênh lệch nhau nhiều nhất là $150$ phút.
}
\end{bt}

% #1965 | Archive/2024/TaiLieuGiangDay/3. K12 - 4in1 - Thongke/data/1. Bai1-Dang1-Khoangbienthien.tex | mã gốc 2D3H1-2
\begin{bt}[Mức độ 2]%[2D3H1-2]
	Thống kê thời gian sử dụng mạng xã hội trong ngày của các bạn Tổ 1, Tổ 2 lớp 12A, được kết quả như sau:
	\begin{center}
		\begin{tabular}{|l|c|c|c|c|}
			\hline
			Thời gian sử dụng (phút) & $[0;10 )$ & $[ 10;30)$ &$[30;60 )$ & $[60;90)$ \\
			\hline
			Số học sinh Tổ 1 & $2$ & $4$ & $3$ & $1$ \\
			\hline
			Số học sinh Tổ 1 & $5$ & $1$ & $3$ & $0$ \\
			\hline
		\end{tabular}
	\end{center}
	Tìm khoảng biến thiên cho thời gian sử dụng mạng xã hội của học sinh mỗi tổ và giải thích ý nghĩa.
	\loigiai{
		Gọi $R_1$, $R_2$ tương ứng là khoảng biến thiên của mẫu số liệu ghép nhóm về thời gian sử dụng mạng xã hội trong ngày của các bạn Tổ $1$ và Tổ $2$.\\
		Ta có: $R_1= 90-0=90$ và $R_2= 60-0=60$.\\
		Do $R_1>R_2$ nên ta có thể kết luận rằng thời gian sử dụng mạng xã hội  trong ngày của các bạn Tổ $1$ phân tán hơn thời gian sử dụng  mạng xã hội của các bạn Tổ $2$.
	}
\end{bt}

% #1966 | Archive/2024/TaiLieuGiangDay/3. K12 - 4in1 - Thongke/data/1. Bai1-Dang1-Khoangbienthien.tex | mã gốc 2D3H1-2
\begin{bt}[Mức độ 2]%[2D3H1-2]
	Dưới đây là kết quả điều tra thời gian hoàn thành bài tập ở nhà môn Toán của 30 học sinh lớp 9:
	\begin{center}
		\textit{	Bảng 3.8. Thời gian hoàn thành bài tập môn Toán ở nhà (đơn vị: phút)}
		\begin{tabular}{|>{\centering\arraybackslash}m{1cm}|*{9}{>{\centering\arraybackslash}m{1cm}|}}
			\hline 
			$42$ & $45$ & $47$ & $47$ & $53$ & $54$ & $58$ & $58$ & $58$ & $59$ \\
			\hline 
			$61$ & $63$ & $64$ & $64$ & $67$ & $68$ & $68$ & $68$ & $70$ & $73$ \\
			\hline 
			$75$ & $75$ & $77$ & $77$ & $78$ & $78$ & $82$ & $82$ & $82$ & $87$ \\
			\hline
		\end{tabular}
	\end{center} Lập mẫu số liệu ghép nhóm với các nhóm ghép có độ dài bằng 10 và nhóm đầu tiên là $[40 ; 50)$. Tìm khoảng biến thiên của mẫu số liệu ghép nhóm.
	\loigiai{Ta có bảng mẫu số liệu ghép nhóm như sau:
		\begin{center}
			\begin{tabular}{|c|c|c|c|}
				\hline Thời gian (phút) &  \begin{tabular}{c} 
					Số học sinh \\
					(Tần số)
				\end{tabular} & \begin{tabular}{c} 
					Tần số \\
					tích luỹ
				\end{tabular} & \begin{tabular}{c} 
					Giá trị \\
					đại diện
				\end{tabular} \\
				\hline$[40 ; 50)$ & $4$ & $4$ & $45$ \\
				\hline$[50 ; 60)$ & $6$ & $10$ & $55$ \\
				\hline$[60 ; 70)$ & $8$ & $18$ & $65$ \\
				\hline$[70 ; 80)$ & $8$ & $26$ & $75$ \\
				\hline$[80 ; 90)$ & $4$ & $30$ & $85$ \\
				\hline
			\end{tabular}	
		\end{center}
		Khoảng biến thiên $R=90-40=50$.
	}
\end{bt}

% #1967 | Archive/2024/TaiLieuGiangDay/3. K12 - 4in1 - Thongke/data/1. Bai1-Dang1-Khoangbienthien.tex | mã gốc 2D3H1-2
\begin{bt}[Mức độ 2]%[2D3H1-2]
	Thời gian trung bình hằng ngày mà một số nhân viên đi từ nhà đến công ty được thống kê trong  \textit{Bảng 3.9}. Tìm khoảng biến thiên của thời gian di chuyển đến công ty của các nhân viên.
	\begin{center}
		\textit{	Bảng 3.9}\\
		\begin{tabular}{|l|c|c|c|c|c|c|}
			\hline Thời gian (phút)  & {$[10 ; 20)$} & {$[20 ; 30)$} & {$[30 ; 40)$} & {$[40 ; 50)$} & {$[50 ; 60)$} & {$[60 ; 70)$} \\
			\hline Số người  & $8$ & $8$ & $10$ & $6$ & $7$ & $1$ \\
			\hline
		\end{tabular}
	\end{center}
	
	\loigiai{
		Khoảng biến thiên $R = 70 - 10 = 60$.
	}
\end{bt}

% #1968 | Archive/2024/TaiLieuGiangDay/3. K12 - 4in1 - Thongke/data/1. Bai1-Dang1-Khoangbienthien.tex | mã gốc 2D3H1-2
\begin{bt}[Mức độ 2]%[2D3H1-2]
Để chuẩn bị mở một trung tâm thể dục thể thao, anh Dũng đã tiến hành điều tra tuổi thọ của máy chạy bộ do hai hãng $X$, $Y$ sản xuất. {\it Bảng 3.4} biểu thị hai mẫu số liệu mà anh thu thập được qua Internet.
\begin{center}
	{\it Bảng 3.4. Tuổi thọ của máy chạy bộ (đơn vị: năm)}
	\begin{tabular}{|p{5cm}|p{1.7cm}|p{1.7cm}|p{1.7cm}|p{1.7cm}|p{1.7cm}|}
		\hline Tuổi thọ &\centering\arraybackslash {$[2; 4)$} &\centering\arraybackslash {$[4; 6)$} &\centering\arraybackslash {$[6; 8)$} &\centering\arraybackslash {$[8; 10)$} &\centering\arraybackslash {$[10; 12)$} \\
		\hline Số máy của hãng $X$&\centering\arraybackslash $7$ &\centering\arraybackslash $20$ &\centering\arraybackslash $36$ &\centering\arraybackslash $20$ &\centering\arraybackslash $17$\\
		\hline Số máy của hãng $Y$&\centering\arraybackslash $0$ &\centering\arraybackslash $20$ &\centering\arraybackslash $35$ &\centering\arraybackslash $35$ &\centering\arraybackslash $10$\\
		\hline
	\end{tabular}
\end{center}
Khoảng biến thiên của mẫu số liệu nào lớn hơn? Từ đó có thể nói là máy chạy bộ do hãng nào sản xuất có tuổi thọ phân tán hơn?
\loigiai{
	Khoảng biến thiên của tuổi thọ máy chạy bộ do hãng $X$ và hãng $Y$ sản xuất tương ứng là $$R_X=12-2=10 \text{ và } R_Y=12-4=8.$$
	Vì $R_X > R_Y$ nên có thể nói là máy do hãng $X$ sản xuất có tuổi thọ phân tán hơn so với máy của hãng $Y$.
} 
\end{bt}

% #1969 | Archive/2024/TaiLieuGiangDay/3. K12 - 4in1 - Thongke/data/1. Bai1-Dang1-Khoangbienthien.tex | mã gốc 2D3H1-2
\begin{bt}[Mức độ 2]%[2D3H1-2]
Người ta tiến hành phỏng vấn hai nhóm khán giả về một bộ phim mới công chiếu. Nhóm $A$ gồm những khán giả thuộc lứa tuổi $20-30$, nhóm $B$ thuộc lứa tuổi trên $30$. Người được hỏi ý kiến phải đánh giá bộ phim bằng cách cho điểm theo một số tiêu chí nêu trong phiếu điều tra và sau đó lấy tổng số điểm (thang điểm $100$). Bảng dưới đây trình bày kết quả điều tra hai nhóm khán giả:
\begin{center}
	{\it Bảng 3.5. Điểm đánh giá của khán giả}\\
	\begin{tabular}{|p{4cm}|p{2cm}|p{2cm}|p{2cm}|p{2cm}|p{2cm}|}
		\hline Điểm &\centering\arraybackslash {$[50; 60)$} &\centering\arraybackslash {$[60; 70)$} &\centering\arraybackslash {$[70; 80)$} &\centering\arraybackslash {$[80;90)$} &\centering\arraybackslash {$[90; 100)$} \\
		\hline Số người của nhóm $A$&\centering\arraybackslash $6$ &\centering\arraybackslash $10$&\centering\arraybackslash $14$ &\centering\arraybackslash $12$ &\centering\arraybackslash $8$ \\
		\hline Số người của nhóm $B$&\centering\arraybackslash $0$ &\centering\arraybackslash $8$ &\centering\arraybackslash $14$ &\centering\arraybackslash $28$ &\centering\arraybackslash $0$ \\
		\hline
	\end{tabular}
\end{center}
Ý kiến đánh giá của nhóm khán giả nào phân tán hơn?
\loigiai{
	Khoảng biến thiên của ý kiến đánh giá của hai nhóm khán giả $A$ và $B$ tương ứng là $$R_A=100-50=50 \text{ và } R_B=90-60=30.$$
	Vì $R_A > R_B$ nên có thể nói là ý kiến đánh giá của nhóm khán giả $A$ phân tán hơn so với ý kiến đánh giá của nhóm khán giả $B$.
}
\end{bt}

% #1970 | Archive/2024/TaiLieuGiangDay/3. K12 - 4in1 - Thongke/data/1. Bai1-Dang1-Khoangbienthien.tex | mã gốc 2D3H1-2
\begin{bt}[Mức độ 2]%[2D3H1-2]
Điểm kiểm tra cuối khoá môn Tiếng Anh của hai lớp ở một trung tâm ngoại ngữ được thống kê trong các {\it Bảng 3.6a} và {\it 3.6b}.

\begin{minipage}{.5\textwidth}
	\begin{center}
		{\it Bảng 3.6a. Điểm của lớp $A$}\\
		\begin{tabular}{|c|c|}
			\hline Điểm & Số học viên (tần số) \\
			\hline $[50; 60)$ & $8$ \\
			\hline $[60; 70)$ & $20$\\
			\hline $[70; 80)$ & $50$\\
			\hline $[80; 90)$ & $17$\\
			\hline $[90; 100)$ & $5$ \\
			\hline
		\end{tabular}
	\end{center}
\end{minipage}
\begin{minipage}{.5\textwidth}
	\begin{center}
		{\it Bảng 3.6b. Điểm của lớp $B$}\\
		\begin{tabular}{|c|c|}
			\hline Điểm & Số học viên (tần số) \\
			\hline $[50; 60)$ & $15$\\
			\hline $[60; 70)$ & $20$\\
			\hline $[70; 80)$ & $30$\\
			\hline $[80; 90)$ & $20$\\
			\hline $[90; 100)$ & $15$\\
			\hline
		\end{tabular}
	\end{center}
\end{minipage}
Tìm khoảng biến thiên của mỗi mẫu số liệu. Có thể dùng khoảng biến thiên để biết điểm của lớp nào đồng đều hơn không?
\loigiai{
	Hai mẫu số liệu đều có khoảng biến thiên là $R=100-50=50$ nên không thể căn cứ vào đó để nói điểm của lớp nào đồng đều hơn.		
}
\end{bt}

% #1971 | Archive/2024/TaiLieuGiangDay/3. K12 - 4in1 - Thongke/data/1. Bai1-Dang1-Khoangbienthien.tex | mã gốc 2D3V1-2
\begin{bt}[Mức độ 3]%[2D3V1-2]
Biểu đồ dưới đây thống kê thời gian tập thể dục buổi sáng mỗi ngày trong tháng 9/2022 của bác Bình và bác An. Ai là người có thời gian tập đều hơn?
\begin{center}
	\begin{tikzpicture}[xscale=2,font=\footnotesize, line join=round, line cap=round, >=stealth]
		\draw[->](0,-0.2)--(0,7)node[left]{\textbf{Số ngày}};
		\draw[->](-0.2,0)--(6,0)node[below]{\textbf{phút}};
		\foreach \i in{1,...,6} \pgfmathsetmacro{\gti}{int(5*(\i))}
		\draw [dotted](0,\i) circle(1pt)node[left]{$\gti$} -- (5.5,\i);
		\foreach \i/\a/\b in {1/15/20,2/20/25,3/25/30,4/30/35,5/35/40}\draw (\i,0) node [below] {[\a;\b)};
		\foreach \i/\tr/\ph in {1/5/0,2/12/25,3/8/5,4/3/0,5/2/0}
		{
			\draw[fill=blue](\i,0)rectangle(\i-0.5,0.2*\tr) node [above right=0.3]{\tr};
			\draw[fill=orange](\i,0)rectangle(\i+0.5,0.2*\ph) node [above left=0.3]{\ph};
		}
		\foreach \i/\a/\b in {3/orange/Bác An, 4/blue/Bác Bình}
		{\draw (6,\i+1)node [right,fill=\a] {};
			\draw (6.2,\i+1) node[right]{\b};
		}
	\end{tikzpicture}
\end{center}
Sử dụng dữ liệu ở biểu đồ trên chọn số thích hợp thay vào các vị trí được đánh dấu ? ở bảng sau
\begin{center}
	\begin{tabular}{|c|c|c|c|c|c|}
		\hline
		Thời gian (phút) & $[15;20)$ & $[20;25)$ & $[25;30)$ & $[30;35)$ & $[35;40)$ \\
		\hline
		Bác Bình & ? & $12$ & $8$ & $3$ & $2$ \\
		\hline
		Bác An & ? & ? & ? & ? & ? \\
		\hline
	\end{tabular}
\end{center}
Hãy tìm khoảng biến thiên của mẫu số liệu ghép nhóm về thời gian tập thể dục buổi sáng mỗi ngày của bác Bình và bác An.
\loigiai{Ta có bảng sau
	\begin{center}
		\begin{tabular}{|c|c|c|c|c|c|}
			\hline
			Thời gian (phút) & $[15;20)$ & $[20;25)$ & $[25;30)$ & $[30;35)$ & $[35;40)$ \\
			\hline
			Bác Bình & $5$ & $12$ & $8$ & $3$ & $2$ \\
			\hline
			Bác An & $0$ & $25$ & $5$ & $0$ & $0$ \\
			\hline
		\end{tabular}
	\end{center}
	Khoảng biến thiên của mẫu số liệu ghép nhóm về thời gian tập thể dục buổi sáng của bác Bình là $R = 40 - 15 = 25$ (phút).\\
	Tuy nhiên, trong mẫu số liệu ghép nhóm về thời gian tập thể dục buổi sáng của bác An, khoảng đầu tiên chứa dữ liệu là $[20; 25)$ và khoảng cuối cùng chứa dữ liệu là $[25; 30)$.\\
	Do đó khoảng biến thiên của mẫu số liệu ghép nhóm về thời gian tập thể dục buổi sáng của bác An là $30 - 20 = 10$ (phút).}
	\end{bt}

% #1972 | Archive/2024/TaiLieuGiangDay/3. K12 - 4in1 - Thongke/data/1. Bai1-Dang1-Khoangbienthien.tex | mã gốc 2D3V1-2
\begin{bt}[Mức độ 3]%[2D3V1-2]
Thầy Tuấn thống kê lại điểm trung bình cuối năm của các học sinh lớp 11A và 11B ở bảng sau:
\begin{center}
	\begin{tabular}{|c|c|c|c|c|c|}
		\hline \diagbox {Lớp}{Điểm trung bình} & {$[5 ; 6)$} & {$[6 ; 7)$} & {$[7 ; 8)$} & {$[8 ; 9)$} & {$[9 ; 10)$} \\
		\hline $11 \mathrm{~A}$ & 1 & 0 & 11 & 22 & 6 \\
		\hline $11 \mathrm{~B}$ & 0 & 6 & 8 & 14 & 12 \\
		\hline
	\end{tabular}
\end{center}
Nếu so sánh theo khoảng biến thiên thì học sinh lớp nào có điểm trung bình ít phân tán hơn?
\loigiai{
	Khoảng biến thiên của điểm số học sinh lớp 11A là: $10-5=5$.\\
	Khoảng biến thiên của điểm số học sinh lớp 11B là: $10-6=4$.\\
	Nếu so sánh theo khoảng biến thiên thì điểm trung bình của các học sinh lớp 11B ít phân tán hơn điểm trung bình của các học sinh lớp 11A.
}
\end{bt}

% #1980 | Archive/2024/TaiLieuGiangDay/3. K12 - 4in1 - Thongke/data/6. Bai2 - Dang2- giatringoaile.tex | mã gốc 2D3H1-3
\begin{vd}[Mức độ 2]%[Bài giảng 12 New-4in1, Phú Thạch]%[2D3H1-3]
			Thời gian (phút) truy cập internet mỗi buổi tối của một số học sinh được cho trong bảng sau:
			\begin{center}
				\begin{tabular}{|c|c|c|c|c|c|c|c|}
					\hline
					Thời gian (phút)&$\left[7{,}5;9{,}5\right)$	& $\left[9{,}5;12{,}5 \right)$ & $\left[12{,}5;15{,}5 \right)$ & $\left[15{,}5;18{,}5 \right)$ & $\left[18{,}5;21{,}5 \right)$ & $\left[21{,}5;24{,}5 \right)$ \\
					\hline
					Số học sinh&$2$&$3$	& $10$ & $15$ &$24$ & $2$ \\
					\hline
				\end{tabular}	
			\end{center}
			Tìm nhóm chứa giá trị ngoại lệ của mẫu số liệu trên.
			\loigiai{
				Cỡ mẫu là $n=2+3+10+15+24+2=56$.\\
				Gọi $x_1$; $x_2$; \ldots; $x_{56}$ là mẫu số liệu đã cho.\\
				Ta có: $x_1$,  $x_{2}\in[7{,5};9{,}5)$; $x_3$, $x_4$, $x_{5}\in [9{,5}; 12{,5})$; $x_{6}$,\ldots $x_{15}\in[12{,}5;15{,}5)$; $x_{16}$, \ldots, $x_{30}\in[15{,}5;18{,}5)$; $x_{31}$, \ldots, $x_{54}\in[18{,}5;21{,}5)$; $x_{55}$, $x_{56}\in[21{,}5;24{,}5)$.\\
				Tứ phân vị thứ nhất $Q_1$ là $\dfrac{x_{14}+x_{15}}{2}$.\\ Do $2$ giá trị $x_{14}$, $x_{15}$ thuộc nhóm $\left[12{,}5;15{,}5 \right)$ nên nhóm này chứa $Q_1$.\\ Do đó, $p=2$; $a_2=12{,}5$; $m_2=12$; $m_1=3$; $a_3-a_2=3$ và ta có $$Q_1=12{,}5+\dfrac{\dfrac{56}{4}-5}{10}\cdot (15{,5}-12{,}5)=15{,}2.$$
				Với tứ phân vị thứ ba $Q_3$ là $\dfrac{x_{42}+x_{43}}{2}$. Do $2$ giá trị $x_{42},\,x_{43}$ thuộc nhóm $\left[18{,}5;21{,}5 \right)$ nên nhóm này chứa $Q_3$.\\ Do đó, $p=4$; $a_4=18{,}5$; $m_4=24$; $m_1+m_2+m_3=3+12+15=30$; $a_5-a_4=3$ và ta có $$Q_3=18{,}5+\dfrac{\dfrac{3\cdot 56}{4}-30}{24}\cdot (21{,}5-18{,}5)=20.$$
				Khoảng tứ phân vị của mẫu số liệu ghép nhóm là $\Delta_Q =Q_3-Q_1=20-15{,}2=4{,}8$.
				Vậy giá $x$ là giá trị ngoại lệ của mẫu số liệu nếu
				$$\hoac{&x<Q_1-1{,}5\Delta_Q\\& x>Q_3+1{,}5\Delta_Q}\Rightarrow \hoac{&x<8\\&x>27{,}2.}$$
				Vậy nhóm chứa giá trị ngoại lệ của mẫu số liệu là $\left[7{,}5;9{,}5\right) $.
			}
		\end{vd}

% #1981 | Archive/2024/TaiLieuGiangDay/3. K12 - 4in1 - Thongke/data/6. Bai2 - Dang2- giatringoaile.tex | mã gốc 2D3H1-3
\begin{vd}[Mức độ 2]%[Bài giảng 12 New-4in1, Phú Thạch]%[2D3H1-3]
			Thời gian chờ khám bệnh của các bệnh nhân tại phòng khám X được cho trong bảng sau:
			\begin{center}
				\begin{tabular}{|l|c|c|c|c|}
					\hline
					Thời gian (phút) & $[0;5 )$ & $[5;10)$ &$[10;15 )$ & $[15;20)$ \\
					\hline
					Số bệnh nhân & $3$ & $12$ & $15$ & $8$ \\
					\hline
				\end{tabular}
			\end{center}
			Mẫu số liệu trên có giá trị ngoại lệ không? Vì sao?
			\loigiai{
				Cỡ mẫu là $n = 3+12+15+8=38$. Gọi $x_1,\ldots,x_{38}$ là thời gian chờ khám bệnh của $38$ bệnh nhân này và giả sử rằng dãy số liệu gốc này đã được sắp xếp theo thứ tự tăng dần.
				Gọi $x_1$; $x_2$; \ldots; $x_{38}$ là mẫu số liệu đã cho.\\
				Ta có: $x_1$,  $x_{2}$, $x_3\in[0;5)$; $x_3$, $x_4$,\ldots, $x_{15}\in [5; 10)$; $x_{16}$,\ldots $x_{30}\in[10;15)$; $x_{31}$, \ldots, $x_{38}\in[15;20)$.\\
				Tứ phân vị thứ nhất của mẫu số liệu gốc là $x_{10}$ nên nhóm chứa tứ phân vị thứ nhất là nhóm $[5;10)$, suy ra
				\[Q_1 = 5+ \dfrac{\left[ \dfrac{38}{4}-3\right] }{12}\cdot (10-5) \approx 7{,}71.\]
				Tứ phân vị thứ ba của mẫu số liệu gốc là $x_{29}$ nên nhóm chứa tứ phân vị thứ ba là nhóm $[10;15)$, suy ra 
				\[Q_3 = 10+ \dfrac{\left[ \dfrac{3\cdot 38}{4}-15\right] }{15}\cdot (15-10) = 14{,}5.\]
				Vậy khoảng tứ phân vị của mẫu số liệu  ghép nhóm là $$\Delta_Q = Q_3 - Q_1 \approx 14{,}5 - 7{,}71= 6{,}79.$$
				Giá trị $x$ là giá trị ngoại lệ của mẫu số liệu nếu
				$$\hoac{&x<Q_1-1{,}5\Delta_Q\\& x>Q_3+1{,}5\Delta_Q}\Rightarrow \hoac{&x<-2{,}475\\&x>24{,}685.}$$	
				Vậy mẫu số liệu trên không có giá trị ngoại lệ.
			}
		\end{vd}

% #1982 | Archive/2024/TaiLieuGiangDay/3. K12 - 4in1 - Thongke/data/6. Bai2 - Dang2- giatringoaile.tex | mã gốc 2D3H1-3
\begin{bt}[Mức độ 2]%[Bài giảng 12 New-4in1, Phú Thạch]%[2D3H1-3]
			Kết quả đo chiều cao của $100$ cây keo 3 năm tuổi tại một nông trường được cho ở bảng sau
			\begin{center}
				\begin{tabular}{|c|c|c|c|c|c|}
					\hline
					Chiều cao (m) & $[8{,}4;8{,}6)$ & $[8{,}6;8{,}8)$ & $[8{,}8;9{,}0)$ & $[9{,}0;9{,}2)$ & $[9{,}2;9{,}4)$ \\
					\hline
					Số cây & $5$ & $12$ & $25$ & $44$ & $14$ \\
					\hline
				\end{tabular}
			\end{center}
			\begin{enumerate}
				\item Hãy tìm khoảng biến thiên và khoảng tứ phân vị của mẫu số liệu ghép nhóm trên.
				\item Trong $100$ cây keo trên có $1$ cây cao $8{,}4$ m. Hỏi chiều cao của cây keo này có phải là giá trị ngoại lệ không?
			\end{enumerate}
			\loigiai{\begin{enumerate}
					\item Khoảng biến thiên của mẫu số liệu là $R=9{,}4-8{,}4=1$.\\
					Ta có cỡ mẫu $n=100$.\\
					Gọi $x_1$; $x_2$; \ldots; $x_{100}$ là mẫu số liệu gồm chiều cao của $100$ cây keo.\\
					Ta có: $x_1$, \ldots, $x_5\in[8{,}4;8{,}6)$; $x_6$, \ldots, $x_{17}\in[8{,}6;8{,}8)$; $x_{18}$, \ldots, $x_{42}\in[8{,}8;9{,}0)$; $x_{43}$, \ldots, $x_{86}\in[9{,}0;9{,}2)$; $x_{87}$, \ldots, $x_{100}\in[9{,}2;9{,}4)$.\\
					Tứ phân vị thứ nhất của mẫu số liệu là $\dfrac{x_{25}+x_{26}}{2}\in[8{,}8;9{,}0)$.\\ Do đó, tứ phân vị thứ nhất của mẫu số liệu ghép nhóm là
					$$Q_1=8{,}8+\dfrac{\dfrac{100}{4}-(5+12)}{25}\cdot(9{,}0-8{,}8)=8{,}864.$$
					Tứ phân vị thứ ba của mẫu số liệu là $\dfrac{x_{75}+x_{76}}{2}\in[9{,}0;9{,}2)$.\\ Do đó, tứ phân vị thứ ba của mẫu số liệu ghép nhóm là
					$$Q_3=9{,}0+\dfrac{\dfrac{3\cdot100}{4}-(5+12+25)}{44}\cdot(9{,}2-9{,}0)=9{,}15.$$
					Vậy khoảng tứ phân vị của mẫu số liệu ghép nhóm là $$\Delta_Q=Q_3-Q_1=0{,}286.$$
					\item Vì $Q_1-1{,}5\Delta_Q=8{,}435$ và $Q_3+1{,}5\Delta_Q=9{,}579$ nên cây keo có chiều cao $8{,}4$ m là giá trị ngoại lệ của mẫu số liệu ghép nhóm.
			\end{enumerate}}
		\end{bt}

% #1983 | Archive/2024/TaiLieuGiangDay/3. K12 - 4in1 - Thongke/data/6. Bai2 - Dang2- giatringoaile.tex | mã gốc 2D3H1-3
\begin{bt}[Mức độ 2]%[Bài giảng 12 New-4in1, Phú Thạch]%[2D3H1-3]
		Kết quả khảo sát năng suất (đơn vị tấn/ha) của một số thửa ruộng được minh họa ở biểu đồ sau
		\begin{center}
			\begin{tikzpicture}[line join=round, line cap=round,>=stealth,xscale=1.5,yscale=1]
				\def\a{1}
				\def\xmax{9}
				\def\ymax{7}
				\tikzset{label style/.style={font=\footnotesize}}
				\draw[->] (0,0)--(\xmax,0) node[below] {\text{Năng suất (tấn/ha)}};
				\draw[->] (0,0)--(0,\ymax) node[below left] {\text{Số thửa ruộng}};
				\draw (0,0) node [below left] {$O$};
				\draw[thin,gray!60]
				(0,1)--(\xmax,1)
				(0,2)--(\xmax,2)
				(0,3)--(\xmax,3)
				(0,4)--(\xmax,4)
				(0,5)--(\xmax,5)
				(0,6)--(\xmax,6)
				;
				\foreach \i in {1,2,3,4,5,6}{
					\draw (0,\i) node[left]{$\i$};
				}
				%		\foreach \i/\j in {0.5*\a/3,1.5*\a/4,2.5*\a/6,3.5*\a/5,4.5*\a/5,5.5*\a/2}{
				%			\draw (\i,\j) node[above]{$\j$};
				%		}
				\foreach \i/\j in {1.5*\a/{[5.5,5.7)},2.5*\a/{[5.7,5.9)},3.5*\a/{[5.9,6.1)},4.5*\a/{[6.1,6.3)},5.5*\a/{[6.3,6.5)},6.5*\a/{[6.5,6.7)}}{
					\draw (\i,-0.5) node[below,rotate=45]{\small $\j$};
				}
				\fill[blue!20]
				(\a,0) rectangle (2*\a,3)
				(2*\a,0) rectangle (3*\a,4)
				(3*\a,0) rectangle (4*\a,6)
				(4*\a,0) rectangle (5*\a,5)
				(5*\a,0) rectangle (6*\a,5)
				(6*\a,0) rectangle (7*\a,2)
				;
				\begin{scope}
					\draw 
					(\a,3)--(2*\a,3) 
					(2*\a,4)--(3*\a,4)
					(3*\a,6)--(4*\a,6)
					(4*\a,5)--(5*\a,5)
					(5*\a,5)--(6*\a,5)
					(6*\a,2)--(7*\a,2)
					(\a,0)--(\a,3)
					(2*\a,0)--(2*\a,4)
					(3*\a,0)--(3*\a,6)
					(4*\a,0)--(4*\a,6)
					(5*\a,0)--(5*\a,5)
					(6*\a,0)--(6*\a,5)
					(7*\a,0)--(7*\a,2)
					;
					\draw (3,\ymax+1) node[above]{\textbf{Năng suất lúa của một số thửa ruộng}};
				\end{scope}
			\end{tikzpicture}
		\end{center}
		\begin{enumerate}
			\item Lập bảng tần số ghép nhóm và tần số tương đối ghép nhóm tương ứng của mẫu số liệu trên.
			\item Hãy xác định nhóm chứa giá trị ngoại lệ của mẫu số liệu trên.
		\end{enumerate}
		\loigiai 
		{
			\begin{enumerate}
				\item Lập bảng tần số ghép nhóm và tần số tương đối ghép nhóm tương ứng của mẫu số liệu trên.
				\begin{center}
					\begin{tabular}{|c|c|c|c|c|c|c|}
						\hline
						Năng suất (tấn/ha) &$[5{,}5;5{,}7)$  & $[5{,}7;5{,}9)$ & $[5{,}9;6{,}1)$ & $[6{,}1;6{,}3)$ & $[6{,}3;6{,}5)$ & $[6{,}5;6{,}7)$\\
						\hline
						Số thửa ruộng	& $3$ & $4$ & $6$ & $5$ & $5$ & $2$\\
						\hline
						Tần số tương đối & $\dfrac{3}{25}$ & $\dfrac{4}{25}$ & $\dfrac{6}{25}$ & $\dfrac{1}{5}$ & $\dfrac{1}{5}$ & $\dfrac{2}{25}$\\
						\hline
					\end{tabular}
				\end{center}
				\item 
			 Cỡ mẫu $n=25$.\\
					Gọi $x_1$; $x_2$;$\ldots$; $x_{25}$ là mẫu số liệu gốc gồm năng suất của $25$ thửa ruộng được sắp xếp theo thứ tự không giảm.\\
					Ta có $x_1$, $x_2$, $x_3\in [5{,}5;5{,}7)$; $x_4$; $\ldots$; $x_7\in [5{,}7;5{,}9)$; $x_8$; $\ldots$; $x_{13}\in [5{,}9;6{,}1)$; \break$x_{14}$; $\ldots$; $x_{18}\in [6{,}1;6{,}3)$; $x_{19}$; $\ldots$; $x_{23}\in [6{,}3;6{,}5)$; $x_{24}$; $x_{25}\in [6{,}5;6{,}7)$.\\
					Tứ phân vị thứ nhất của mẫu số liệu gốc là $\dfrac{1}{2}(x_6+x_7)\in [5{,}7;5{,}9)$. Do đó, tứ phân vị thứ nhất của mẫu số liệu ghép nhóm là 
					$$Q_1=5{,}7+\dfrac{\dfrac{25}{4}-3}{4}(5{,}9-5{,}7)=5{,}8625.$$
					Tứ phân vị thứ ba của mẫu số liệu gốc là $\dfrac{1}{2}(x_{19}+x_{20})\in [6{,}3;6{,}5)$. Do đó, tứ phân vị thứ ba của mẫu số liệu ghép nhóm là 
					$$Q_3=6{,}3+\dfrac{\dfrac{3\cdot 25}{4}-(3+4+6+5)}{5}(6{,}5-6{,}3)=6{,}33.$$
					Khoảng tứ phân vị của mẫu số liệu ghép nhóm là 
					$$\Delta_Q=6{,}33-5{,}8625=0{,}4675.$$		Giá trị $x$ là giá trị ngoại lệ của mẫu số liệu nếu
					$$\hoac{&x<Q_1-1{,}5\Delta_Q\\& x>Q_3+1{,}5\Delta_Q}\Rightarrow \hoac{&x<5{,}16125\\&x>7{,}03125.}$$
			Vậy mẫu số liệu trên không chứa giá trị ngoại lệ.					
			\end{enumerate}
		}
	\end{bt}

% #1984 | Archive/2024/TaiLieuGiangDay/3. K12 - 4in1 - Thongke/data/6. Bai2 - Dang2- giatringoaile.tex | mã gốc 2D3V1-3
\begin{bt}[Mức độ 3]%[Bài giảng 12 New-4in1, Phú Thạch]%[2D3V1-3]
		Hằng ngày ông Thắng đều đi xe buýt từ nhà đến cơ quan. Dưới đây là bảng thống kê thời gian của $100$ lần ông Thắng đi xe buýt từ nhà đến cơ quan.
		\begin{center}
			\begin{tabular}{|c|c|c|c|c|c|c|}
				\hline
				Thời gian (phút) & $[15;18)$ & $[18;21)$ & $[21;24)$ & $[24;27)$ & $[27;30)$ & $[30;33)$ \\
				\hline
				Số lượt & $22$ & $38$ & $27$ & $8$ & $4$ & $1$ \\
				\hline
			\end{tabular}
		\end{center}
		\begin{enumerate}
			\item Hãy tìm khoảng tứ phân vị của mẫu số liệu ghép nhóm trên.
			\item Biết rằng trong $100$ lần đi trên, chỉ có đúng một lần ông Thắng đi hết hơn $29$ phút. Thời gian của lần đi đó có phải là giá trị ngoại lệ không?
			\item 		Hãy tìm khoảng biến thiên và khoảng tứ phân vị của mẫu số liệu ghép nhóm ở trên sau khi đã loại bỏ các giá trị ngoại lệ. Em có nhận xét gì về khoảng biến thiên, khoảng tứ phân vị vừa tìm được và khoảng biến thiên, khoảng tứ phân vị ban đầu?
		\end{enumerate}
		\loigiai{\begin{enumerate}
				\item Cỡ mẫu $n=100$.\\
				Gọi $x_1$; $x_2$; \ldots; $x_{100}$ là mẫu số liệu gốc gồm thời gian $100$ lần đi xe buýt của ông Thắng.\\
				Ta có: $x_1$, \ldots, $x_{22}\in[15;18)$; $x_{23}$, \ldots, $x_{60}\in[18;21)$; $x_{61}$, \ldots, $x_{87}\in[21;24)$; $x_{88}$, \ldots, $x_{95}\in[24;27)$; $x_{96}$, \ldots, $x_{99}\in[27;30)$; $x_{100}\in[30;33)$.\\
				Tứ phân vị thứ nhất của mẫu số liệu gốc là $\dfrac{x_{25}+x_{26}}{2}\in[18;21)$. \\Do đó, tứ phân vị thứ nhất của mẫu số liệu ghép nhóm là
				$$Q_1=18+\dfrac{\dfrac{100}{4}-22}{38}\cdot(21-18)=\dfrac{693}{38}.$$
				Tứ phân vị thứ ba của mẫu số liệu gốc là $\dfrac{x_{75}+x_{76}}{2}\in[21;24)$.\\ Do đó, tứ phân vị thứ ba của mẫu số liệu ghép nhóm là
				$$Q_3=21+\dfrac{\dfrac{3\cdot100}{4}-(22+38)}{27}\cdot(24-21)=\dfrac{68}{3}.$$
				Vậy khoảng tứ phân vị của mẫu số liệu ghép nhóm là $$\Delta_Q=Q_3-Q_1=\dfrac{505}{114}.$$
				\item Trong lần duy nhất ông Thắng đi hết hơn $29$ phút, thời gian đi của ông thuộc nhóm $[30;33)$.\\
				Vì $Q_3+1{,}5\Delta_Q=\dfrac{6\,683}{228}<30$ nên thời gian của lần ông Thắng đi hết hơn $29$ phút là giá trị ngoại lệ của mẫu số liệu ghép nhóm.
				\item 
				Ta có $Q_3+1{,}5\Delta_Q=\dfrac{6683}{228}<30$ và $Q_1-1{,}5\Delta_Q=\dfrac{881}{76}\approx11{,}59$.\\
				Trong lần duy nhất ông Thắng đi hết hơn $29$ phút, thời gian đi của ông thuộc nhóm $[30;33)$, do đó các giá trị thuộc nhóm $[30;33)$ là các giá trị ngoại lệ, sau khi loại bỏ các giá trị ngoại lệ, ta có bảng thống kê sau đây:
				\begin{center}
					\begin{tabular}{|c|c|c|c|c|c|}
						\hline
						Thời gian (phút) & $[15;18)$ & $[18;21)$ & $[21;24)$ & $[24;27)$ & $[27;30)$ \\
						\hline
						Số lượt & $22$ & $38$ & $27$ & $8$ & $4$ \\
						\hline
					\end{tabular}
				\end{center}
				Khoảng biến thiên của mẫu số liệu khi chưa loại bỏ các giá trị ngoại lệ là $R=33-15=18$.\\
				Khoảng biến thiên của mẫu số liệu sau khi loại bỏ các giá trị ngoại lệ là $R=30-15=15$.\\
				Sau khi loại bỏ các giá trị ngoại lệ, mẫu số liệu mới có độ phân tán thấp hơn.\\
				Xét về khoảng tứ phân vị, ta có cỡ mẫu $n=99$.\\
				Gọi $x_1$; $x_2$; \ldots; $x_{99}$ là mẫu số liệu gốc gồm thời gian $99$ lần đi xe buýt của ông Thắng.\\ Ta có: $x_1$, \ldots, $x_{22}\in[15;18)$; $x_{23}$, \ldots, $x_{60}\in[18;21)$; $x_{61}$, \ldots, $x_{87}\in[21;24)$; $x_{88}$, \ldots, $x_{95}\in[24;27)$; $x_{96}$, \ldots, $x_{99}\in[27;30)$.\\
				Tứ phân vị thứ nhất của mẫu số liệu gốc là $x_{25}\in[18;21)$.\\ Do đó, tứ phân vị thứ nhất của mẫu số liệu ghép nhóm là
				$$Q_1=18+\dfrac{\dfrac{99}{4}-22}{38}\cdot(21-18)=\dfrac{2769}{152}.$$
				Tứ phân vị thứ ba của mẫu số liệu gốc là $x_{75}\in[21;24)$.\\ Do đó, tứ phân vị thứ ba của mẫu số liệu ghép nhóm là
				$$Q_3=21+\dfrac{\dfrac{3\cdot99}{4}-(22+38)}{27}\cdot(24-21)=\dfrac{271}{12}.$$
				Vậy khoảng tứ phân vị của mẫu số liệu ghép nhóm sau khi loại bỏ các giá trị ngoại lệ vẫn là $$\Delta_Q=Q_3-Q_1=\dfrac{1991}{456}.$$
				Vậy khoảng tứ phân vị của mẫu số liệu sau khi loại bỏ giá trị ngoại lệ nhỏ hơn khoảng tứ phân vị ban đầu.
		\end{enumerate}}
	\end{bt}
